\documentclass[10pt,a4paper]{article}
\usepackage{amsmath,amssymb,amsthm,graphicx}
\usepackage{array,aliascnt,ifthen,hyperref,cleveref,enumerate,xcolor,mathtools}
\newtheorem{theorem}{Theorem}[section]
\newtheorem{lemma}[theorem]{Lemma}
\newtheorem{proposition}[theorem]{Proposition}
\newtheorem{corollary}[theorem]{Corollary}
\theoremstyle{definition}
\newtheorem{definition}[theorem]{Definition}
\newtheorem{example}[theorem]{Example}
\theoremstyle{remark}
\newtheorem{remark}[theorem]{Remark}
\AtBeginDocument{\special{pdf:mapfile fdsymbol.map}}

\numberwithin{equation}{section}


 

\usepackage{ifthen}

\newboolean{spellingmode}
\setboolean{spellingmode}{false}

\newboolean{proofsinappendix}
\setboolean{proofsinappendix}{false}




\usepackage{xspace}
\usepackage{fdsymbol}
\usepackage{tikz}
\usepackage{ifdraft}
\usepackage{booktabs}
\usepackage{environ}

\usepackage{hyperref}
% Debug-only showkeys omitted in r002 presentation.
\usepackage{microtype}





\ifthenelse{\boolean{spellingmode}}{	
	\AtEndPreamble{
	\hyphenpenalty=50000 		\pagestyle{empty}	\thispagestyle{empty}	\def\pagestyle#1{}	\def\thispagestyle#1{}		\def\labelmarginpar#1{}	}
}{}

\ifdraft{
\overfullrule=4pt }{
\overfullrule=0pt }



\newcommand{\overbar}[1]{\mkern 1.5mu\overline{\mkern-1.5mu#1\mkern-1.5mu}\mkern 1.5mu}
\newcommand{\mybar}[3]{  \mathrlap{\hspace{#2}\overline{\scalebox{#1}[1]{\phantom{\ensuremath{#3}}}}}\ensuremath{#3}
}

\newcommand{\myhat}[3]{  \mathrlap{\hspace{#2}\widehat{\scalebox{#1}[1]{\phantom{\ensuremath{#3}}}}}\ensuremath{#3}
}

\newcommand{\barF}{\mybar{0.6}{2.5pt}{F}}
\newcommand{\barH}{\mybar{0.6}{2.5pt}{H}}
\newcommand{\barB}{\mybar{0.6}{2.5pt}{B}}
\newcommand{\barBs}{\mybar{0.6}{1.8pt}{B}}
\newcommand{\barPow}{\mybar{0.7}{1.5pt}{\Pow}}
\newcommand{\barSigmas}{\mybar{0.9}{0pt}{\Sigmas}}
\newcommand{\barrho}{\mybar{0.9}{1pt}{\rho}}

\newcommand{\gra}{\mathsf{gra}}
\newcommand{\dom}{\mathsf{dom}}

\DeclareMathOperator{\Sub}{\mathrm{Sub}}


\providecommand{\catname}{\mathsf} 
\providecommand{\clsname}{\mathcal}
\providecommand{\oname}[1]{{\mathop{\mathsf{#1}}\xspace}}

\def\defcatname#1{\expandafter\def\csname B#1\endcsname{\catname{#1}}}
\def\defcatnames#1{\ifx#1\defcatnames\else\defcatname#1\expandafter\defcatnames\fi}
\defcatnames ABCDEFGHIJKLMNOPQRSTUVWXYZ\defcatnames

\def\defclsname#1{\expandafter\def\csname C#1\endcsname{\clsname{#1}}}
\def\defclsnames#1{\ifx#1\defclsnames\else\defclsname#1\expandafter\defclsnames\fi}
\defclsnames ABCDEFGHIJKLMNOPQRSTUVWXYZ\defclsnames


\def\Set{\catname{Set}}
\def\Pos{\catname{Pos}}
\def\Nom{\catname{Nom}}
\def\Vect{\catname{Vec}_K}
\def\Conv{\catname{Conv}}
\def\Cpo{\catname{Cpo}}
\def\Cppo{\catname{Cppo}}
\def\dCpo{\catname{dCpo}}
\def\Top{\catname{Top}}

\newcommand{\Konig}{Kőnig}

\providecommand{\eps}{{\operatorname\epsilon}}
\providecommand{\amp}{\mathbin{\&}}
\providecommand{\from}{\leftarrow}

\providecommand{\argument}{-}
\providecommand{\altargument}{\underline{\;\;}}

\providecommand{\wave}[1]{\widetilde{#1}}		               \providecommand{\ul}{\underline}			                     
\DeclareOldFontCommand{\bf}{\normalfont\bfseries}{\mathbf}
\providecommand{\mplus}{{\scriptscriptstyle\bf+}} 	       
\providecommand{\PSet}{{\mathcal P}}			                 \providecommand{\FSet}{{\mathcal P}_{\omega}}		           \providecommand{\CSet}{{\mathcal P}_{\omega_1}}		         \providecommand{\NESet}{{\mathcal P}^{\mplus}}		         \providecommand{\Id}{\operatorname{Id}}

\providecommand{\Hom}{\mathsf{Hom}}
\providecommand{\id}{\mathsf{id}}
\providecommand{\op}{\mathsf{op}}
\providecommand{\comp}{\mathbin{\circ}}
\providecommand{\iso}{\mathbin{\cong}}
\providecommand{\tensor}{\mathbin{\otimes}}
\providecommand{\unit}{\star}				
\providecommand{\bang}{\operatorname!}				             
\providecommand{\ito}{\hookrightarrow}				             \providecommand{\pto}{\mathrel{\rightharpoonup}}		       \providecommand{\tto}{\mathrel{\Rightarrow}}			         \providecommand{\mto}{\mapsto}
\providecommand{\xto}[1]{\,\xrightarrow{#1}\,}
\providecommand{\xfrom}[1]{\,\xleftarrow{\;#1}\,}
\providecommand{\To}{\mathrel{\Rightarrow}}			           \providecommand{\From}{\mathrel{\Leftarrow}}			         \providecommand{\flatxto}[1]{\xto{#1}}


\providecommand{\dar}{\kern-1.2pt\operatorname{\downarrow}}	
\providecommand{\uar}{\kern-1.2pt\operatorname{\uparrow}}	

\providecommand{\True}{\top}
\providecommand{\False}{\bot}

\providecommand{\bigor}{\bigvee}
\providecommand{\bigand}{\bigwedge}

\providecommand{\impl}{\Rightarrow}
\providecommand{\equ}{\Longleftrightarrow}

\providecommand{\entails}{\vdash}

\providecommand{\appr}{\sqsubseteq}
\providecommand{\join}{\sqcup}
\providecommand{\meet}{\sqcap}

\providecommand{\bigjoin}{\bigsqcup}
\providecommand{\bigmeet}{\bigsqcap}

\providecommand{\fst}{\oname{fst}}
\providecommand{\snd}{\oname{snd}}
\providecommand{\pr}{\prime}

\providecommand{\brks}[1]{\langle #1\rangle}
\providecommand{\Brks}[1]{\bigl\langle #1\bigr\rangle}

\providecommand{\inl}{\oname{inl}}
\providecommand{\inr}{\oname{inr}}
\providecommand{\inj}{\oname{in}}

\DeclareSymbolFont{Symbols}{OMS}{cmsy}{m}{n}
\DeclareMathSymbol{\iobj}{\mathord}{Symbols}{"3B}

\providecommand{\curry}{\oname{curry}}
\providecommand{\uncurry}{\oname{uncurry}}
\providecommand{\ev}{\oname{ev}}

\usepackage{stmaryrd}

\providecommand{\lsem}{\llbracket}
\providecommand{\rsem}{\rrbracket}
\providecommand{\sem}[1]{\lsem #1 \rsem}

\providecommand{\Lsem}{\bigl\llbracket}
\providecommand{\Rsem}{\bigr\rrbracket}
\providecommand{\Sem}[1]{\Lsem #1 \Rsem}

\providecommand{\comma}{,\operatorname{}\linebreak[1]}		 \providecommand{\dash}{\nobreakdash-\hspace{0pt}}		       \providecommand{\erule}{\rule{0pt}{0pt}}		                                                                          
\providecommand{\by}[1]{\text{/\!\!/~#1}}			             \providecommand{\pacman}[1]{}					                     
\newcommand{\undefine}[1]{\let #1\relax}					                       
\providecommand{\noqed}{\def\qed{}}				                 
\providecommand{\mone}{{\text{\kern.5pt\rmfamily-}\mathsf{\kern-.5pt1}}}

\providecommand{\smin}{\smallsetminus}




 
\makeatletter
\def\mfix#1{\oname{#1}\@ifnextchar\bgroup\@mfix{}}	       \def\@mfix#1{#1\@ifnextchar\bgroup\mfix{}}			           \makeatother

\providecommand{\ift}[3]{\mfix{if}{\mathbin{}#1}{then}{\mathbin{}#2}{else}{\mathbin{}#3}}
\providecommand{\case}[3]{\mfix{case}{\mathbin{}#1}{of}{#2}{\kern-1pt;}{\mathbin{}#3}}




\usepackage{dsfont}

\DeclareMathSymbol{\mathinvertedexclamationmark}{\mathord}{operators}{'074}
\DeclareMathSymbol{\mathexclamationmark}{\mathord}{operators}{'041}
\makeatletter
\newcommand{\raisedmathinvertedexclamationmark}{  \mathord{\mathpalette\raised@mathinvertedexclamationmark\relax}}
\newcommand{\raised@mathinvertedexclamationmark}[2]{  \raisebox{\depth}{$\m@th#1\mathinvertedexclamationmark$}}
\makeatother

\newcommand{\sqdown}{\logp{\Downarrow}}

\newcommand{\lfplong}{locally finitely presentable\xspace}
\newcommand{\Lfplong}{Locally finitely presentable\xspace}
\newcommand{\lfp}{\lfplong}
\newcommand{\Lfp}{\Lfplong}

\newcommand{\strength}{\mathsf{st}}
\newcommand{\product}{\times}
\newcommand{\wt}{\widetilde}
\newcommand{\longTo}{\Longrightarrow}
\newcommand{\Pt}{V}
\newcommand{\R}{\mathcal{R}}
\newcommand{\lsmc}{\fatsemi}
\newcommand{\Pre}{\mathbf{Pre}}
\newcommand{\Pow}{\mathcal{P}}
\newcommand{\xTo}{\xRightarrow}
\newcommand{\HJBis}{\mathbf{HJBisim}}
\newcommand{\Bis}{\mathbf{Bisim}}
\newcommand{\primrec}{\mathsf{pr}\,}
\newcommand{\hattau}{\hat{\tau}}
\newcommand{\hatini}{\hat{\ini}}
\newcommand{\hatr}{\hat{r}}
\newcommand{\hatbeta}{\hat{\beta}}
\newcommand{\lft}{{\mathsf{l}}}
\newcommand{\rgt}{{\mathsf{r}}}
\DeclareMathOperator{\colim}{\mathsf{colim}}
\newcommand{\terminal}{\mathexclamationmark}
\newcommand{\initial}{\raisedmathinvertedexclamationmark}
\newcommand{\smc}{\mathbin{;}}
\newcommand{\hatR}{{\widehat{R}}}
\newcommand{\qand}{\quad\text{and}\quad}
\newcommand{\qqand}{\qquad\text{and}\qquad}
\newcommand{\rsto}{\rightsquigarrow}
\newcommand{\pb}[2]{\scalebox{#1}{\rotatebox{#2}{$\ulcorner$}}}
\newcommand{\deq}[1]{{\rotatebox{#1}{$=$}}}
\newcommand{\dleq}[1]{{\rotatebox{#1}{$\preceq$}}}
\newcommand{\dgeq}[1]{{\rotatebox{#1}{$\succeq$}}}
\newcommand{\cev}[1]{\reflectbox{\ensuremath{\vec{\reflectbox{\ensuremath{#1}}}}}}
\newcommand{\diag}{\mathsf{in}}
\newcommand{\HO}{\mathcal{HO}}
\newcommand{\Refl}{\mathbf{Refl}}
\newcommand{\Gra}{\mathbf{Gra}}
\newcommand{\RGra}{\mathbf{RGra}}
\newcommand{\RRel}{\mathbf{RRel}}
\newcommand{\happrox}{\mathrel{\widehat\lesssim}}
\newcommand{\happroxs}{\mathrel{\hat\approx^{\star}}}
\newcommand{\under}[1]{\lvert#1\rvert}
\newcommand{\SKI}{\mathrm{SKI}}
\renewcommand{\L}{\mathcal{L}}
\newcommand{\ap}{{\mathrm{ap}}}
\newcommand{\var}{\mathsf{var}}
\newcommand{\cn}{{\mathrm{cn}}}
\newcommand{\cv}{{\mathrm{cv}}}
\newcommand{\Powf}{\mathcal{P}_{\omega}}
\newcommand{\FSigma}{{\mathbb{\Sigma}}}
\newcommand{\FSigmas}{{\mathbb{\Sigma}}^\star}
\newcommand{\Sigmas}{\Sigma^{\star}}
\newcommand{\FDelta}{{\mathbb{\Delta}}}
\newcommand{\FDeltas}{{\mathbb{\Delta}}^\star}
\newcommand{\ar}{\mathsf{ar}}
\newcommand{\NT}{\mathrm{Nat}}
\newcommand{\supp}{\mathsf{supp}}
\newcommand{\epito}{\twoheadrightarrow}
\newcommand{\abs}{\mathsf{abs}}
\newcommand{\conc}{\mathsf{conc}}
\newcommand{\T}{{\mathcal{T}}}
\renewcommand{\S}{{\mathcal{S}}}
\newcommand{\can}{\mathsf{can}}
\newcommand{\lab}{\mathsf{label}}
\newcommand{\leaves}{\mathsf{leaves}}
\newcommand{\incl}{\mathsf{in}}
\newcommand{\seq}{\subseteq}
\newcommand{\ol}{\overline}
\newcommand{\out}{\mathsf{out}}
\newcommand{\outl}{\mathsf{l}}
\newcommand{\outr}{\mathsf{r}}
\newcommand{\trc}{{\mathsf{trc}}}
\newcommand{\fn}{{\mathsf{fn}}}
\newcommand{\beh}{{\mathsf{beh}}}
\newcommand{\dbtilde}[1]{\widetilde{\raisebox{0pt}[0.85\height]{$\widetilde{#1}$}}}
\providecommand{\C}{}
\providecommand{\D}{}
\providecommand{\E}{}

\newcommand{\fp}{{\mathrm{fp}}}
\newcommand{\fg}{{\mathrm{fg}}}
\newcommand{\rec}{{\mathrm{rec}}}
\newcommand{\wf}{{\mathrm{wf}}}
\newcommand{\frc}{\Coalg_{\fp,\rec}(H)}
\newcommand{\fwf}{\Coalg_{\fp,\wf}(H)}
\newcommand{\fgwf}{\Coalg_{\fp,\wf}(H)}


\renewcommand{\C}{{\mathbb{C}}}
\renewcommand{\D}{{\mathbb{D}}}
\renewcommand{\E}{{\mathbb{E}}}
\newcommand{\I}{{\mathbb{I}}}

\newcommand{\A}{\mathcal{A}}

\newcommand{\iref}[2]{\autoref{#1}.\ref{#1:#2}}


\renewcommand{\id}{{\mathsf{id}}}
\newcommand{\Nat}{\mathds{N}}
\renewcommand{\NT}{\mathsf{Nat}}
\newcommand{\Rel}{\mathbf{Rel}}

\newcommand{\ngoesv}[2]{\ensuremath{#1 \centernot{\xrightarrow{~#2~}}}} 
\newcommand{\Goes}[2]{\ensuremath{#1 \Rightarrow #2}}
\newcommand{\goes}[2]{\ensuremath{#1 \rightarrow #2}}
\newcommand{\Rets}[2]{\ensuremath{#1 \Downarrow #2}}
\newcommand{\Retsc}[2]{\ensuremath{#1 \Rightarrow #2, \checkmark}}
 \newcommand{\rets}[2]{\ensuremath{#1 \downarrow #2}}
\newcommand{\retsc}[2]{\ensuremath{#1 \rightarrow #2, \checkmark}}

\newcommand{\Goesv}[3]{\ensuremath{#1 \xRightarrow{~#3~} #2}}
\newcommand{\goesv}[3]{\ensuremath{#1 \xrightarrow{~#3~} #2}}
\newcommand{\retsv}[3]{\ensuremath{#1 \xrightarrow{~#3~} #2, \checkmark}}

\newcommand{\retsvr}[3]{\ensuremath{#1~\textcolor{red}{\Downarrow_{#3}}~#2}}


\newcommand{\stimes}[1]{{S\product #1}}
\newcommand{\xtimes}[1]{{X\product #1}}
\newcommand{\ytimes}[1]{{Y\product #1}}
\newcommand{\sbrks}[1]{{ \brks{\fst,#1}  }}
\newcommand{\cpair}[2]{{ #1, #2 }}
\newcommand{\tpair}[2]{{ #1 \product #2 }}


\renewcommand{\brack}[1]{\lbrack #1 \rbrack}
\newcommand{\sbrack}[1]{\llbracket #1 \rrbracket}
\newcommand{\ssbrack}[1]{\llbracket\mkern-4.7mu\llbracket #1 \rrbracket\mkern-4.7mu\rrbracket}
\newcommand{\f}{\oname{f}}
\newcommand{\fii}{\oname{f}}

\newcommand{\takeout}[1]{\empty}

\newcommand{\Dist}{\mathcal{D}}
\newcommand{\CPowf}{\Pow_{\mathsf{c},\omega}}

\newcommand{\ini}{\iota}
\newcommand{\ter}{\tau}
\newcommand{\wh}{\widehat}
\DeclareMathOperator{\Coalg}{\mathsf{Coalg}}
\DeclareMathOperator{\Alg}{\mathsf{Alg}}
\newcommand{\cofun}{(\overline{\rule{-1pt}{7pt}\argument})}
\renewcommand{\rho}{\varrho}
\newcommand{\N}{\mathds{N}}
\newcommand{\opp}{\mathsf{op}}
\newcommand{\xla}[1]{\xleftarrow{~#1~}}
\newcommand{\mypowfin}{\mathscr{P}_\omega}

\newcommand{\pullbackangle}[2][]{\arrow[phantom,to path={
                     -- ($ (\tikztostart)!1cm!#2:([xshift=8cm]\tikztostart) $)
                        node[anchor=west,pos=0.0,rotate=#2,
                        inner xsep = 0]
                        {\begin{tikzpicture}[minimum
                        height=1mm,baseline=0,#1]
    \draw[-] (0,0) -- (.5em,.5em) -- (0,1em);
                        \end{tikzpicture}}}]{}}

\makeatletter
\newsavebox{\@brx}
\newcommand{\llangle}[1][]{\savebox{\@brx}{\(\m@th{#1\langle}\)}  \mathopen{\copy\@brx\kern-0.5\wd\@brx\usebox{\@brx}}}
\newcommand{\rrangle}[1][]{\savebox{\@brx}{\(\m@th{#1\rangle}\)}  \mathclose{\copy\@brx\kern-0.5\wd\@brx\usebox{\@brx}}}
\makeatother

\renewcommand{\comp}{\cdot}
\renewcommand{\c}{\colon}

\newcommand{\Var}{\mathop{\mathsf{Var}}}
\newcommand{\xra}{\xrightarrow}
\renewcommand{\xto}{\xra}
\let\xmpsto=\xmapsto

\newcommand{\V}{\mathcal{V}}
\newcommand{\monoto}{\rightarrowtail}
\newcommand{\subto}{\hookrightarrow}
\newcommand{\commu}{\ensuremath{\circlearrowleft}}
\renewcommand{\Nat}{\mathbb{N}}
\newcommand{\pfun}{\rightharpoonup}
\newcommand{\fset}{{\mathbb{F}}}
\newcommand{\vcat}{{\Set}^{\fset}}
\newcommand{\alg}[1]{\mathbf{Alg}(#1)}
\newcommand{\coalg}[1]{\mathbf{CoAlg}(#1)}
\newcommand{\mon}{\bullet}

\newcommand{\act}[1]{\mathbf{Act}(#1)}
\newcommand{\adj}[4]{#1\negmedspace: #2\rightleftarrows #3:\negmedspace #4}
\newcommand{\gcat}{\mathbb{C}}
\providecommand{\Dinat}{\mathsf{Dinat}}
\providecommand{\DiNat}{\mathsf{DiNat}}
\newcommand{\mS}{{\mu\Sigma}}
\DeclareMathOperator{\iter}{\oname{it}}
\DeclareMathOperator{\coiter}{\oname{coit}}
\newcommand{\transpose}[1]{\overline{#1}}
\renewcommand{\epsilon}{\varepsilon}

\newcommand*\xbar[1]{  \kern.2em\hbox{    \vbox{      \hrule height 0.5pt       \kern0.5ex      \hbox{        \kern-0.2em        \ensuremath{#1}        \kern-0.4em      }    }  }\kern.4em } 

\makeatletter
\newcommand{\monto}{\@ifstar{\@mtolifted}{\@mto}}
\newcommand{\@mto}{\multimapdot}
\newcommand{\@mtolifted}{\mathbin{\xbar{\multimapdot}}}
\makeatother

\newcommand{\Ar}[1]{\mathbf{Ar}(#1)}

\newcommand{\bool}{\mathsf{bool}}
\newcommand{\true}{\mathsf{true}}
\newcommand{\false}{\mathsf{false}}
\newcommand{\ite}[3]{\mathsf{if}\,{#1}\,\mathsf{then}\,#2\,\mathsf{else}\,#3}
\newcommand{\app}[2]{\,}
\newcommand{\skiterm}{\Lambda_{\varsigma}}
\newcommand{\skitermu}{\Lambda_{u}}
\newcommand{\finalc}{Z}



\usetikzlibrary{cd,calc}

\def\resettheorembrackets{
\def\theorembracketopen{(}
\def\theorembracketclose{)}
}
\resettheorembrackets
\makeatletter
\def\@spopargbegintheorem#1#2#3#4#5{\trivlist
      \item[\hskip\labelsep{#4#1\ #2}]{#4{\theorembracketopen}#3{\theorembracketclose}\@thmcounterend\ }#5}
\makeatother


\newcommand{\resetCurThmBraces}{  \gdef\curThmBraceOpen{(}  \gdef\curThmBraceClose{)}}
\resetCurThmBraces
\newcommand{\removeThmBraces}{  \gdef\curThmBraceOpen{}  \gdef\curThmBraceClose{}}
\resetCurThmBraces

\newenvironment{notheorembrackets}{\removeThmBraces}{\resetCurThmBraces}

\makeatletter




\newcommand{\proofappendixbegin}[2]{  \phantomsection  \subsection*{\textbf{#1~\autoref{#2}}}  \addcontentsline{toc}{subsection}{#1~\autoref{#2}}  \label{#2:proof}  \def\proofappendix@qedsymbolmissing{\qed}
}
\newcommand{\proofappendixend}{  \proofappendix@qedsymbolmissing}
\let\oldqedhere\qedhere
\def\qedhere{\global\def\proofappendix@qedsymbolmissing{}\oldqedhere}

\newenvironment{proofhere}[2][Proof of]{    \proof{#1~\autoref{#2}}    }{  \endproof  \par}

    \newenvironment{proofappendix}[2][Proof of]{    \begin{proofhere}[#1]{#2}\label{#2:proof}
  }{    \end{proofhere}
  }

\makeatother



\tikzstyle{shiftarr}=[
        rounded corners,        to path={--([#1]\tikztostart.center)
                     -- ([#1]\tikztotarget.center) \tikztonodes
                     -- (\tikztotarget)},
]

\tikzset{
    commutative diagrams/.cd,
    arrow style=tikz,
    diagrams={>={Straight Barb[length=1.75pt,width=3.85pt,inset=1.95pt]}},     row sep=large,
    column sep = huge
}

\tikzset{cong/.style={draw=none,edge node={node [sloped, allow upside down, auto=false]{$\cong$}}},
         iso/.style={draw=none,every to/.append style={edge node={node [sloped, allow upside down, auto=false]{$\cong$}}}}}


\usetikzlibrary{decorations.pathmorphing}

\tikzcdset{scale cd/.style={every label/.append style={scale=#1},
    cells={nodes={scale=#1}}}}

\newcommand{\set}[2][]{  \ifthenelse{\equal{#2}{}}{    \ensuremath{{#1\emptyset}}  }{    \ensuremath{{#1\{#2#1\}}}  }}

\newcommand{\Z}{\ensuremath{\mathbb{Z}}}

\newcommand{\simple}{\ensuremath{\overset{0}{\rightarrowtail}}\xspace}
\newcommand{\injmono}{\ensuremath{\overset{\smash{\raisebox{-1pt}{\ensuremath{\scriptstyle+}}}}{\rightarrowtail}}}\newcommand{\hpresint}{\ensuremath{H\cap}\xspace}



\usepackage{etoolbox}
\patchcmd{\thmhead}{(#3)}{\curThmBraceOpen #3\curThmBraceClose }{}{}

\theoremstyle{plain}

\theoremstyle{definition}
\newcommand{\lipicsnewtheorem}[2]{  \newaliascnt{#1}{theorem}
  \newtheorem{#1}[#1]{#2}
  \aliascntresetthe{#1}
  \expandafter\newcommand\csname #1autorefname\endcsname{#2}
}
\lipicsnewtheorem{notation}{Notation}
\lipicsnewtheorem{construction}{Construction}
\lipicsnewtheorem{assumption}{Assumption}
\lipicsnewtheorem{rem}{Remark}

\AtBeginDocument{
\def\subsectionautorefname{Section}
}

\newcommand{\textqt}[1]{`#1'}
\newcommand{\Perm}{\mathfrak{S}}
\newcommand{\fs}{\mathsf{fs}}
\newcommand{\of}{\mathsf{of}}
\newcommand{\At}{\mathbb{A}}
\bibliographystyle{plainurl}
\title{Well-Founded Coalgebras Meet \Konig's Lemma} 
\patchcmd{\author}{Anonymous author}{Anonymous author(s)}{}{}

\makeatletter
\ifx\@hideLIPIcs\@undefined
\else
\renewcommand\keywordsHeading{\def\@keywords{\vspace{-1.8\baselineskip}}}
\renewcommand\subjclassHeading{\def\@ccsdescString{\vspace{-1.8\baselineskip}}}
\fi
\makeatother


\providecommand{\twnote}[2][]{#2}
\title{Coalgebraic Konig lemma for well-founded coalgebras over locally finitely presentable categories}
\author{Henning Urbat, Thorsten Wißmann}
\date{arxiv:2507.18539v3; CC BY4.0 (fixed full author version)}
\begin{document}
\maketitle
\noindent\textit{Editorial note.} Independent CC BY4.0 fixed full author version with complete original proofappendix bodies.\par\medskip
\noindent\textit{Editorial note.} Only main Coalgebraic Konig lemma counted. Supporting specialized proofs are in proof spans and not split into inventory.\par\medskip
\noindent\textit{Editorial note.} Original proofsinappendix=false branch prints identical author proofs inline; publishing/final unused note infrastructure omitted from wrapper only.\par\medskip
\noindent\textit{Editorial note.} Native FdSymbol style/fonts/map and full source/licenses bundled locally; no shared installation or local compilation.\par\medskip
\noindent\textit{Editorial note.} Original FXRegisterAuthor creates twnote. Actual complete body has one twnote[inline]\{\} call with empty content. Wrapper restores an optional-argument twnote that prints every nonempty content argument verbatim and suppresses only emptiness; original body/source unchanged.\par\medskip
\noindent\textit{Editorial note.} r002 presentation only: root visual inspection found original showkeys debug boxes covering proof prose. New independent preamble-r002 removes showkeys; successful r001 preamble/job, native originals and all mathematical body remain unchanged. Same counted theorem; no extra item.\par\medskip
\noindent\textit{Editorial scope.} Complete author material collection for one structural theorem; author retains responsibility for the mathematics.\par\medskip
\noindent\textit{Allowed standard prerequisites (editorial).}\begin{itemize}
\item Locally finitely presentable categories, finitely generated objects, strong epi/mono factorization and finitary functors. Exact original Assumption3.1 and all named mono/intersection conditions retained. Standard category theory facts are explicitly cited to Adamek/Rosicky and original references.
\item Well-founded coalgebra background. Complete author definitions/preliminaries and original references retained. Actual new coproduct coalgebra construction and its full well-foundedness proof are retained as core support, not reclassified as prerequisite.
\end{itemize}\medskip
\section{Introduction}
In 1927, Dénes \Konig~\cite{konig1927} proved a result in graph theory, known as \emph{\Konig's lemma}, that turned out to be of fundamental importance. In its most popular formulation, the lemma asserts that every infinite, finitely branching tree contains an infinite path. While this statement looks almost obvious (with a classical, non-constructive mindset at least), it isolates an important selection principle that appears in different shapes and forms throughout mathematics and computer science, for instance, in proofs of compactness and completeness theorems in logic~\cite{teller89}, Ramsey theory~\cite{szekely23}, or the analysis of infinite runs and liveness properties of systems in verification~\cite{brv04} and automata theory~\cite{gmms14,isb16}. \Konig's lemma is moreover deeply tied with the foundations of set theory, specifically the axiom of countable choice~\cite{hr98,jech02}. Variants such as the \emph{weak \Konig's lemma} (the restriction to binary trees) and Brouwer's \emph{fan theorem}~\cite{brouwer27} have been extensively studied in (reverse) constructive mathematics and computability theory~\cite{di21,hirschfeldt14}.

\Konig's lemma applies to graphs, which are the simplest instance of \emph{state-based systems}, that is, systems with a notion of \emph{state space} and a notion of \emph{transition} between states. Such systems can be uniformly modelled at categorical generality using \emph{coalgebras}~\cite{jacobs16,rutten00}. A {coalgebra} is given by an endofunctor $H$ on a category $\C$ that specifies the {system type}, an object $C$ of $\C$ determining the state space, and a morphism $c\colon C\to HC$ determining the transition structure. By varying the base category $\C$ and the type functor $H$, countless different incarnations of state-based systems emerge in the coalgebraic framework. For example, graphs are precisely coalgebras for the power set functor $\Pow$ on the category of sets, non-deterministic register automata~\cite{KaminskiFrancez94} over an infinite alphabet $\mathds{A}$ of data values correspond to coalgebras for the functor $\{0,1\}\times \Pow^{\mathds{A}}$ on the topos of nominal sets~\cite{BojanczykEA14} (with power object functor $\Pow$), and (determinized) probabilistic automata are coalgebras for a functor that involves the convex power set functor on the category of convex sets~\cite{bss17}.

In the present work, we aim to address the following natural question:
\[ \emph{Which coalgebras (beyond graphs) feature a \Konig's lemma?} \]
To motivate our approach, let us revisit the case of graphs ($\Pow$-coalgebras). Note first that finitely branching graphs are precisely coalgebras for the \emph{finite} power set functor $\Powf$. We say that a graph is \emph{well-founded} if it contains no infinite paths. Then König's lemma can be restated, in contrapositive form and in coalgebraic terminology, as follows:
\begin{equation}\label{eq:konig-pow} \emph{Every state of a well-founded $\Powf$-coalgebra lies in some finite subcoalgebra}. \end{equation}
The initial observation towards a generalization of \eqref{eq:konig-pow} is that the notion of well-foundedness is available for \emph{arbitrary} coalgebras~\cite{taylor99,amm25}: A coalgebra $c\colon C\to HC$ is \emph{well-founded} if it has no proper {cartesian} subcoalgebras (see \autoref{sec:prelim} for details). Moreover, $\Powf$ is the prototypical \emph{finitary} functor, that is, a functor that preserves filtered colimits. The first coalgebraic extension of \Konig's lemma, recently proved by Ad\'amek, Milius, and Moss~\cite[Prop.~9.2.19]{amm25}, replaces the finite power set functor $\Powf$ with an arbitrary finitary set functor $H$:
\begin{equation}\label{eq:konig-fin-set} \emph{Every state of a well-founded $H$-coalgebra lies in some finite subcoalgebra}. \end{equation}
The proof from \emph{op.~cit.}\ exploits the fact that every coalgebra $c\colon C\to HC$ for a finitary set functor can be associated with a \emph{canonical graph} $\ol{c}\colon C\to \Powf C$ representing the transition structure of the coalgebra. In this way, \eqref{eq:konig-fin-set} reduces to the classical \Konig's lemma \eqref{eq:konig-pow}. 

The main contribution of our paper is a substantial generalization of \eqref{eq:konig-fin-set} to coalgebras over \emph{abstract} categories. This step involves two non-trivial challenges: First, canonical graphs are specific to coalgebras over the category of sets, so the reduction argument by Ad\'amek~et~al.~\cite{amm25} no longer works. Second, it is not clear what a `finite' coalgebra is in a general category. To address these challenges, we base our theory on \emph{\lfp} categories (which include, for instance, the categories of sets, nominal sets, and convex sets). The key feature of \lfplong categories is that they come with two well-behaved abstract notions of `finite' object, namely \emph{finitely presentable} and \emph{finitely generated} objects. Finite coalgebras thus naturally generalize to coalgebras with either finitely presentable or finitely generated state space, or \emph{fp/fg-carried coalgebras} for short. Our main result (\autoref{thm:koenig}) establishes \Konig's lemma at this level of generality:

\medskip\noindent \textsf{\bfseries Coalgebraic \Konig's Lemma.} \emph{Let $\C$ be a \lfplong category with monic coproduct injections, and let $H$ be a finitary functor on $\C$ preserving binary intersections. Then every well-founded $H$-coalgebra is the directed join of its fg-carried well-founded subcoalgebras.}

\medskip\noindent In \autoref{sec:applications}, we explore several instantiations of this result beyond the category of sets. Notably, we demonstrate that every \lfplong topos features a \Konig's lemma for graphs, and we derive versions of \Konig's lemma for nominal and convex transition systems. The latter systems are \emph{infinitely} branching, but their branching is finitely generated in an algebraic sense, which turns out to be enough to guarantee existence of infinite paths.

The core ingredient for our proof of the Coalgebraic \Konig's Lemma is a general construction on coalgebras, \emph{coproduct extension}~\cite{amv06,urbat17,lffiac,mpw16}, which allows to add new states to a coalgebra in a restricted way. The key observation is that coproduct extension preserves well-foundedness. This fact leads us to a further fundamental result on well-founded coalgebras, namely a new construction of the \emph{initial algebra} for the functor $H$ (\autoref{thm:wf-initial-algebra}):

\medskip\noindent \textsf{\bfseries Initial Algebra Theorem.} \emph{Under the assumptions of the Coalgebraic \Konig's Lemma, the initial algebra for the functor $H$ is the colimit of all fp-carried well-founded coalgebras.}

\medskip\noindent This theorem relates to a recent result by Wißmann and Milius~\cite{wm24}, who showed that the initial algebra is the colimit of all fp-carried \emph{recursive} coalgebras~\cite{taylor99}. Under the assumptions of our Initial Algebra Theorem, well-founded coalgebras form a (in some instances proper) subclass of recursive ones. Thus,
the Initial Algebra Theorem can be seen as an improvement of the corresponding result for recursive coalgebras because it operates with a smaller diagram, and more importantly, well-foundedness of a given coalgebra is typically much easier to prove than recursivity in settings where the two notions do not coincide. In addition, the idea underlying the proof of the Initial Algebra Theorem carries over to its recursive version and yields a transparent new proof of the latter that greatly simplifies the original argument~\cite{wm24}.


\section{Preliminaries}\label{sec:prelim}
We start with some background on \lfplong categories~\cite{adamek_rosicky_1994} and on well-founded and recursive coalgebras~\cite{amm25,taylor99}; see the cited textbooks for more details. Readers should be familiar with basic category theory~\cite{mac2013categories}, such as functors and (co)limits.



\subparagraph*{Locally Finitely Presentable Categories}
The results of our paper apply to coalgebras over \lfplong categories, which are categories where every object can be approximated from below by suitably `finite' objects. For example, the category of sets is \lfplong because every set is a union of finite sets, and the category of vector spaces is \lfplong because every vector space is a  union of finite-dimensional spaces. The categorical notion of finiteness is expressed abstractly using filtered colimits.

A category $I$ is \emph{filtered} if (i) $I$ is non-empty, (ii) for any two objects $i,j\in I$ there exists a cospan $i\to k \leftarrow j$ in $I$, and (iii) for every parallel pair $f,g\colon i\to j$ there exists a morphism $h\colon j\to k$ such that $h\circ f = h\circ g$. For example, finitely cocomplete categories are filtered. A poset (viewed as a category) is filtered iff it is directed, i.e.\ every finite subset has an upper bound. A \emph{filtered diagram} $D\colon I\to \C$ is a diagram whose scheme $I$ is filtered, and a \emph{filtered colimit} is a colimit of a filtered diagram. A functor is \emph{finitary} if it preserves filtered colimits.

\begin{example}[Filtered Colimits in $\Set$]\label{ex:filtered-set}
In the category $\Set$ of sets and functions, the colimit of a filtered diagram $D\colon I\to \Set$ is given by $(\coprod_i D_i)/{\approx}$ where $\approx$ is the equivalence relation with  $(x,i)\approx (y,j)$ iff there exist morphisms $f\colon i\to k$ and $g\colon j\to k$ in $I$ such that $Df(x)=Dg(y)$. The colimit injection $c_i\colon D_i\to (\coprod_i D_i)/{\approx}$ sends $x\in D_i$ to the equivalence class of $(x,i)$. An endofunctor $H$ on $\Set$ is finitary iff for every set $X$ and every $x\in HX$, there exists a finite subset $m\colon M\subto X$ such that $x\in Hm[HM]$~\cite[Rem.\ 3.14]{amsw19}.
\end{example}

An object $X$ of a category $\C$ is \emph{finitely presentable} if its hom-functor $\C(X,-)\colon \C\to \Set$ is finitary. More explicitly, this means that given a filtered diagram $D\colon I\to \C$ with colimit cocone $(c_i\colon D_i\to C)_{i\in I}$, every morphism from $X$ to $C$ \emph{factorizes essentially uniquely through the cocone $(c_i)$}, that is, the following two statements hold:
\begin{enumerate}
\item every morphism $f\colon X\to C$ factorizes as $f=c_i\circ g$ for some $i\in I$ and $g\colon X\to D_i$;
\item for any two morphisms $g,g'\colon X\to D_i$ with $i\in I$ and $c_i\circ g = c_i\circ g'$, there exists $k\colon i\to j$ in $I$ such that $Dk\circ g = Dk\circ g'$.
\end{enumerate}
Besides finitely presentable objects, there is also the weaker notion of \emph{finitely generated} object. An object $X$ of $\C$ is \emph{finitely generated} if for every filtered diagram $D\colon I\to \C$ whose colimit cocone $(c_i\colon D_i\monoto C)_{i\in I}$ consists of monomorphisms, every morphism from $X$ to $C$ factorizes through some $c_i$. (Here the essential uniqueness of the factorization comes for free because $c_i$ is monic.) We let $\C_\fp,\C_\fg\hookrightarrow \C$ denote the full subcategories given by finitely presentable and finitely generated objects, respectively. Note that $\C_\fp\seq \C_\fg$, but generally finitely presentable and finitely generated objects do not coincide. 

\begin{notheorembrackets}
\begin{lemma}\label{lem:fp-fp-fin-colimits} The subcategories $\C_\fp$, $\C_\fg \subto \C$ have the following closure properties:
\begin{enumerate}
\item Both $\C_\fp$ and $\C_\fg$ are closed under finite colimits~\cite[Prop.~1.3]{adamek_rosicky_1994}. 
\item $\C_\fg$ is closed under strong quotients (represented by strong epimorphisms)~\cite[Prop.~1.69]{adamek_rosicky_1994}.
\end{enumerate}
\end{lemma}
\end{notheorembrackets}

A category $\C$ is \emph{\lfplong} if it is cocomplete, and there exists a (small) set~$\A$ of finitely presentable objects such that every object of $\C$ is a filtered colimit of objects in $\A$, that is, a colimit of some filtered diagram $D\colon I\to \C$ such that $D_i\in \A$ for all $i\in I$. Informally, this says that every object can be constructed from `finite' objects.

\begin{example}[Locally Finitely Presentable Categories]\label{ex:categories} 
\begin{figure*}[t]
\centering
\def\arraystretch{1.0}
\begin{tabular}{@{}p{1.5cm}p{2.5cm}p{3.2cm}p{5.5cm}@{}}
\toprule
$\C$ & Objects & Morphisms & $\C_\fp$ ($=\C_\fg$) \\
\midrule
$\Set$ & sets & functions & finite sets~ \\
$\Nom$ & nominal sets &  equivariant maps & orbit-finite sets \\
$\Conv$ & convex sets & affine maps & finitely generated convex sets \\
\bottomrule
\end{tabular}
\caption{\Lfplong categories.}\label{fig:categories}
\end{figure*}
In our applications (\autoref{sec:applications}) we consider the \lfp categories of \autoref{fig:categories}. Their finitely presentable and finitely generated objects coincide and are described in the last column. The categories $\Nom$ of nominal sets and $\Conv$ of convex sets are discussed in more detail in \autoref{sec:nom} and \autoref{sec:convex}. Let us mention that there are many more examples of \lfp categories, including all categories of algebras for an equational theory (such as monoids, rings, vector spaces), all categories of relational structures axiomatized by a finitary relational Horn theory (such as posets), and all presheaf categories $\Set^\D$ over a small category $\D$~\cite{adamek_rosicky_1994}. In all these categories, filtered colimits are formed like in $\Set$ (\autoref{ex:filtered-set}).
\end{example}

\begin{rem}\label{rem:loc-pres}
There is a more general notion of \emph{locally $\lambda$-presentable category}, for a regular cardinal $\lambda$, where the role of filtered colimits is taken over by \emph{$\lambda$-filtered colimits}. Here a diagram scheme is \emph{$\lambda$-filtered} if condition (ii) in the above definition of a filtered category is strengthened to ($\text{ii}_\lambda$): Every non-empty set of objects of cardinality less than $\lambda$ has a cospan over it. \Lfp categories correspond to the case $\lambda=\omega$. A category is \emph{locally presentable} if it is locally $\lambda$-presentable for some $\lambda$. An example of a locally ($\omega_1$-)presentable category that is not locally \emph{finitely} presentable is the category of metric spaces and non-expansive maps. The results of our paper currently do not extend to arbitrary locally presentable categories; in fact, we use several results on well-founded coalgebras (see below) which rest on properties specific to ($\omega$-)filtered colimits.
\end{rem}




\Lfp categories are, in many respects, convenient and well-behaved structures. Let us mention a few core results, whose proofs can be found in Ad\'amek and Rosick\'y~\cite{adamek_rosicky_1994} (Remark 1.56, Proposition 1.61, Proposition 1.62, Exercise~1.o).

\begin{lemma}\label{lem:lfp-complete-wp}
Every \lfp category $\C$ is complete and well-powered, and every morphism in $\C$ has a (strong epi, mono)-factorization.
\end{lemma}
Since a \lfp category $\C$ is cocomplete and has (strong epi, mono)-factorizations, the poset $\Sub(X)$ of subobjects of an object $X$ form a complete lattice: The join of a family of subobjects $m_i\colon X_i\monoto X$ ($i\in I$) is the subobject $m\colon \bigvee_{i\in I} X_i \monoto X$ given by the (strong epi, mono)-factorization of the morphism $[m_i]_{i\in I}$:
\[ 
\begin{tikzcd}[outer sep=0pt]
\coprod_{i\in I} X_i
  \ar[shiftarr={yshift=15}]{rr}{[m_i]_{i\in I}}
  \ar[two heads]{r}
  & \bigvee_{i\in I} X_i \ar[tail]{r}{m}
  & X
   \end{tikzcd}
\]
\emph{Directed} joins of subobjects can be alternatively computed by just forming their colimit. In more detail, a directed family $m_i\colon X_i\monoto X$ ($i\in I$) of subobjects can be regarded as a cocone over the directed diagram $D\colon I\to \C$, $i\mapsto X_i$, where $I$ is ordered by $i\leq j$ iff $m_i\leq m_j$ in the usual partial order of subobjects. Then the following statement holds:


\begin{lemma}\label{lem:lfp-directed-unions}
Let $\C$ be \lfp. For any directed family $m_i\colon X_i\monoto X$ ($i\in I$) of subobjects of an object $X$, its join in $\Sub(X)$ is the colimit of the induced diagram.
\end{lemma}

We also have a useful characterization of filtered colimits in general:

\begin{lemma}\label{lem:filtered-colimits}
Let $D\colon I\to \C$ be a filtered diagram in a \lfp category~$\C$. A cocone $(c_i\colon D_i\to C)_{i\in I}$ over $D$ forms a colimit cocone iff for every $X\in \C_\fp$, every morphism from $X$ to $C$ factorizes essentially uniquely through $(c_i)$.
\end{lemma}
Note that the left-to-right implication holds by the definition of a finitely presentable object. The reverse implication is the interesting part of the lemma. 

Regarding finitary functors, the following result generalizing \autoref{ex:filtered-set} applies:

\begin{notheorembrackets}
\begin{lemma}[{\cite[Thm.~3.4]{amsw19}}]\label{lem:finitary-criterion}
Let $\C$ be a \lfp category with $\C_\fp = \C_\fg$, and let $H$ be a functor that preserves monomorphisms. Then $H$ is finitary iff for every object~$X$ of $\C$, every finitely generated subobject $s\colon S\monoto HX$ factorizes through $Hm\colon HM\monoto HX$ for some finitely generated subobject $m\colon M\monoto X$.  
\end{lemma}
\end{notheorembrackets}
This result applies, in particular, to mono-preserving functors on the categories of \autoref{fig:categories}.

Lastly, if a functor is not finitary, it can be canonically restricted to a finitary functor. In the following, let $[\C,\C]$ denote the category of endofunctors on $\C$ and natural transformations, $[\C,\C]_{\mathsf{fin}}$ its full subcategory given by finitary endofunctors, and $\C_\fp/X$ the full subcategory of the slice category $\C/X$ whose objects are morphisms $(f\colon P\to X)$ of $\C$ with $P\in \C_\fp$.

\begin{notheorembrackets}
\begin{lemma}[{\cite[Cor.~2.8]{amv03}}]\label{lem:finitary-coreflection}
If $\C$ is \lfp, then $[\C,\C]_{\mathsf{fin}}$ is a coreflective subcategory of $[\C,\C]$. The coreflection of $H\colon \C\to \C$ is the functor $H_\omega\colon \C\to\C$ given by $H_\omega X=\colim D_X$ for the diagram $D_X\colon \C_\fp/X\to \C$ mapping $(f\colon P\to X)\in \C_\fp/X$ to $HP$.   
\end{lemma}
\end{notheorembrackets}
The functor $H_\omega$ is called the \emph{finitary coreflection} of $H$. For example, the finitary coreflection of the power set functor $\Pow\colon \Set\to \Set$ is the finite power set functor $\Powf\colon \Set\to \Set$.

\subparagraph*{Coalgebras and Algebras} 

Coalgebras~\cite{jacobs16,rutten00} form a uniform categorical abstraction of state-based transition systems, such as graphs, labelled transition systems, Markov chains, and various kinds of automata. The idea is to model the state space of a system by an object of a suitable category $\C$, and its transition type by an endofunctor on that category. Formally, a \emph{coalgebra} for an endofunctor $H\colon \C\to \C$ is a pair $(C,c)$ consisting of an object $C$ (the
\emph{state space}) and a morphism $c\colon C\to HC$ (its
\emph{structure}). A \emph{morphism} $h\colon (C,c)\to (D,d)$ of coalgebras is a morphism $h\colon C\to D$ of $\C$ such that $Hh\circ c = d\circ h$. We let $\Coalg(H)$ denote the category of coalgebras for $H$ and their morphisms. 

\emph{Subcoalgebras} of a coalgebra $(C,c)$ are represented by coalgebra morphisms  $m\colon (S,s)\monoto (C,c)$ with $m$ monic in $\C$, and \emph{strong quotients} of $(C,c)$ are represented by coalgebra morphisms $e\colon (C,c)\epito (Q,q)$ with $e$ strongly epic in $\C$. If $H$ preserves monomorphisms, a subcoalgebra is uniquely determined by the object $S$. If moreover $\C$ has (strong epi, mono)-factorizations, then every coalgebra morphism $h\colon (C,c)\to (D,d)$ factorizes uniquely as a strong quotient followed by a subcoalgebra: $h = (\begin{tikzcd} (C,c) \ar[two heads]{r}{e} & (S,s) \ar[tail]{r}{m} & (D,d) \end{tikzcd})$. This factorization is called the \emph{image factorization} of $h$, and $m$ is the \emph{image} of $h$.

\begin{example}[Graphs]\label{ex:coalgebras}
 A coalgebra for the power set functor $\Pow$ on $\Set$ is precisely a (directed) graph. A coalgebra for the finite power set functor $\Powf$ (the finitary coreflection of $\Pow$) corresponds to a \emph{finitely branching} graph, i.e.\ one where every node has only finitely many outgoing edges. Given a graph $(C,c)$, we write $x\to y$ iff $y\in c(x)$. 
A subcoalgebra of $(C,c)$ is a subset $S\seq C$ closed under successors: $x\in S$ and $x\to y$ implies $y\in S$.
\end{example}

\begin{example}[Automata and Transition Systems]
Various important types of automata and transition systems can be modelled as coalgebras for a set functor $H$, including deterministic automata ($HX=\{0,1\}\times X^A$ for a fixed finite input alphabet $A$), non-deterministic automata ($HX=\{0,1\}\times (\Pow X)^A$), labelled transition systems ($HX = (\Pow X)^A$), discrete labelled Markov chains ($HX=(\mathcal{D}X)^A$ where $\mathcal{D}$ is the finite distribution functor), and many more~\cite{jacobs16,rutten00}.
\end{example}

Examples of coalgebras beyond the category of sets are discussed later, e.g.\ graphs in a topos (\autoref{sec:topos}), coalgebras in nominal sets which model automata and transition systems over
  infinite alphabets (\autoref{sec:nom}), and 
  coalgebras in convex sets, modelling systems that combine probabilistic and non-deterministic branching
  (\autoref{sec:convex}).

We refer the reader to Jacobs~\cite{jacobs16} or Rutten~\cite{rutten00} for general introductions to the rich theory of coalgebras. For our purposes, we only need to recall two basic results. The first one asserts that colimits in $\Coalg(H)$ are formed in the underlying category:

\begin{notheorembrackets}
\begin{lemma}[{\cite[Prop.~4.7]{rutten00}}]\label{lem:colimits-coalg}
The forgetful functor from $\Coalg(H)$ to $\C$ creates colimits.
\end{lemma}
\end{notheorembrackets}

The second result concerns coalgebras with `finite' state space. A coalgebra is \emph{fp-carried} if its state space is finitely presentable, and \emph{fg-carried} if its state space is finitely generated. We let $\Coalg_\fp(H), \Coalg_\fg(H)\subto \Coalg(H)$ denote the corresponding full subcategories. 

\begin{notheorembrackets}
\begin{lemma}[{\cite[Lem.~3.2]{ap04}}]\label{lem:fp-fg-coalg} Let $H\colon \C\to \C$ be a finitary endofunctor.
\begin{enumerate}
\item Every $fp$-carried coalgebra is a finitely presentable object of $\Coalg(H)$.
\item If $H$ preserves monos, every fg-carried coalgebra is a finitely generated object of $\Coalg(H)$.
\end{enumerate}
\end{lemma}
\end{notheorembrackets}

Dually to the notion of coalgebra, \emph{algebras} for an endofunctor yield a categorical abstraction of algebraic structures. An \emph{algebra} $(A,a)$ for $H\colon \C\to \C$ is given by an object~$A$ (its \emph{carrier}) and a morphism $a\colon H A\to A$ (its \emph{structure}). A
\emph{morphism} $h\colon (A,a) \to (B,b)$ of algebras is a morphism
$h\colon A\to B$ of~$\gcat$ such that $h\circ a = b\circ H h$.  Algebras
for $H$ and their morphisms form a category $\Alg(H)$. An \emph{initial algebra} $(I,i)$ is an initial object of $\Alg(H)$. By Lambek's lemma~\cite{lambek68}, its structure $i\colon HI\to I$ is an isomorphism in $\C$.



\begin{example}[$\Sigma$-Algebras]
The prime example of functor algebras are algebras for a
signature. An \emph{(algebraic) signature} is given by a set~$\Sigma$
of \emph{operation symbols} and a map $\ar\colon \Sigma\to \Nat$
associating to every $\f\in \Sigma$ its \emph{arity}. Every
signature~$\Sigma$ induces the polynomial endofunctor $H_\Sigma X = \coprod_{\f\in \Sigma} X^{\ar(\f)}$ on $\Set$. An algebra for $H_\Sigma$ is
precisely an algebra for the signature~$\Sigma$: a set $A$ with an operation $\f^A\colon A^n\to A$ for every $n$-ary operation symbol $\f\in \Sigma$. The initial algebra for $H_\Sigma$ is the algebra of closed terms formed over the signature $\Sigma$.
\end{example}




\subparagraph*{Well-Founded and Recursive Coalgebras}\label{sec:prelim:wf-rec}
The concept of `having no infinite paths' can be extended from graphs to arbitrary coalgebras~\cite{taylor99}. A subcoalgebra $m\colon (S,s)\monoto (C,c)$ of a coalgebra $(C,c)$ is \emph{cartesian} if the commutative diagram in \autoref{fig:cartesian} is a pullback. A coalgebra $(C,c)$ is \emph{well-founded} if it is has no proper cartesian subcoalgebras: for every cartesian subcoalgebra $m\colon (S,s)\monoto (C,c)$, the monomorphism $m$ is an isomorphism.
\begin{figure}
  \begin{minipage}{.45\textwidth}\centering
  \begin{tikzcd}[row sep=5mm]
  S \pullbackangle{-45} \ar{r}{s} \ar[tail]{d}[swap]{m}  & HS \ar{d}{Hm} \\
  C \ar{r}{c} & HC 
  \end{tikzcd}
  \caption{A cartesian subcoalgebra.}
  \label{fig:cartesian}
  \end{minipage}\hfill  \begin{minipage}{.5\textwidth}\centering
\begin{tikzcd}[row sep=5mm]
C \ar{r}{h} \ar{d}[swap]{c} & A \\
HC \ar{r}{Hh} & HA \ar{u}[swap]{a}
\end{tikzcd}
  \caption{A coalgebra-to-algebra morphism.}
  \label{fig:coalg2alg}
  \end{minipage}
\end{figure}



\begin{example}[Graphs]\label{ex:wfCoalg}
Let $(C,c)$ be a $\Pow$-coalgebra, viewed as a graph. We note first that a subset $S\seq C$ carries a cartesian subcoalgebra iff, for all $x\in C$,
\begin{equation}\label{eq:cartesian-pow}
x\in S \iff \text{all successors of $x$ lie in $S$.}
\end{equation}
Indeed, the left-to-right implication says that the square in \autoref{fig:cartesian} commutes ($S$ is a subcoalgebra), and the right-to-left implication says that it is a pullback. From this, it follows that
\[ \text{$(C,c)$ is well-founded} \iff \text{$(C,c)$ has no infinite paths}. \]
 For the left-to-right implication, suppose that $(C,c)$ is well-founded. Then the set $S\seq C$ of all states that lie on no infinite path is a cartesian subcoalgebra by \eqref{eq:cartesian-pow}, and so $S=C$. For the right-to-left implication, we argue contrapositively. Suppose that $(C,c)$ is not well-founded, and let $S\subsetneq C$ be a proper cartesian subcoalgebra. Pick $x_0\in C\setminus S$ arbitrarily. By \eqref{eq:cartesian-pow}, some successor $x_1$ of $x_0$ lies in $C\setminus S$. By \eqref{eq:cartesian-pow} again, some successor $x_2$ of $x_1$ lies in $C\setminus S$. Repeating this argument yields an infinite path $x_0\to x_1\to x_2\to \cdots$.
\end{example}

\begin{example}[Coalgebras for Set Functors]
For every functor $H\colon \Set\to \Set$ that preserves wide intersections and every set $X$, there is map $\supp_X\colon HX\to \Pow X$ that sends an element $t\in HX$ to its \emph{support}, the least subset $m\colon M\subto X$ such that $t\in Hm[HM]$. Every coalgebra $(C,c)$ for $H$ thus induces a coalgebra $(C,\supp_C\circ c)$ for~$\Pow$, the \emph{canonical graph} of $(C,c)$. A coalgebra for $H$ is well-founded iff its canonical graph is well-founded~\cite[Rem.~6.3.4]{taylor99}.
\end{example}
One important application of well-founded graphs is the principle of \emph{well-founded induction} and \emph{well-founded recursion}. This principle can be understood at the level of coalgebras~\cite{taylor99}. 

Given a coalgebra $(C,c)$ and an algebra $(A,a)$ for the same functor $H$, a \emph{coalgebra-to-algebra morphism} from $(C,c)$ to $(A,a)$ is a morphism $h\colon C\to A$ making the square in \autoref{fig:coalg2alg} commute. A coalgebra $(C,c)$ is \emph{recursive} if for every algebra $(A,a)$ there exists a unique coalgebra-to-algebra morphism from $(C,c)$ to $(A,a)$.

\begin{example}[Graphs]
Every well-founded $\Pow$-coalgebra $(C,c)$ is recursive. Indeed, this amounts to saying that for every $a\colon \Pow A \to A$, there is unique map $h\colon C\to A$ such that $h(x)=a(h[c(x)])$ for all $x\in C$. This is precisely the principle of \emph{well-founded recursion}. The special case where $a\colon \Pow\{0,1\}\to \{0,1\}$ sends $\emptyset$, $\{1\}$ to $1$ and $\{0\}$, $\{0,1\}$ to $0$ corresponds to \emph{well-founded induction}: given a predicate $P\colon C\to \{0,1\}$, one has
\[ \big(\,\forall x\in C.\, (\forall y\in c(x).\, P(y) \implies P(x))\,\big) \quad\implies\quad \big(\,\forall x\in C.\, P(x)\,\big). \] 
\end{example}
\noindent
In the following, we denote by
\[\Coalg_{\fp,\wf}(H),\, \Coalg_{\fg,\wf}(H),\, \Coalg_\wf(H),\, \Coalg_{\fp,\rec}(H),\, \Coalg_{\fg,\rec}(H),\, \Coalg_\rec(H) \] 
the full subcategories of $\Coalg(H)$ given by all (fp-carried, fg-carried)  well-founded coalgebras, and analogously for recursive coalgebras. 

We conclude with some general properties of well-founded and recursive coalgebras. First, colimits of recursive coalgebras are computed in the underlying category~\cite[Prop.~7.1.9]{amm25}:
\begin{lemma}\label{lem:rec-coalg-colimits}
The subcategory $\Coalg_\rec(H)\subto \Coalg(H)$ is closed under colimits.
\end{lemma}
We also mention that recursive coalgebras bear a close connection with initial algebras: If $(I,i)$ is an initial algebra for $H$, then $(I,i^{-1})$ is a final recursive coalgebra, and vice versa~\cite{cuv06}.

Regarding well-founded coalgebras, we need the results of the proposition below. We state them for a \lfp base category $\C$, our setting of interest; however, they hold under slightly weaker assumptions~\cite[Cor.~8.4.3, Thm.~8.5.3]{amm25}. In the following, an initial object $0$ is called \emph{simple} if for every object $X$ the unique morphism $0\to X$ is monic.


\begin{proposition}\label{prop:wf} 
Let $\C$ be a \lfp category with a simple initial object, and let $H\colon \C\to \C$ be an endofunctor that preserves monomorphisms.
\begin{enumerate}
\item\label{prop:wf:closed-colim-quotients} The subcategory $\Coalg_\wf(H)\subto \Coalg(H)$ is closed under colimits and strong quotients.
\item\label{prop:wf:wf-to-rec} Every well-founded coalgebra is recursive, that is, $\Coalg_\wf(H)\seq \Coalg_\rec(H)$.
\end{enumerate}
\end{proposition}








\section{\Konig's Lemma for Coalgebras}\label{sec:konig}
We are prepared to work towards the core result of our paper, \Konig's lemma for coalgebras over a \lfp category (\autoref{thm:koenig}). Let us first fix the global setup:

\begin{assumption}\label{asm}
For the remainder of this paper, we fix a finitary endofunctor $H$ on a \lfplong category $\C$. 
\end{assumption}
Most of our results require the following two additional conditions on $\C$ and $H$:

\medskip\noindent (\injmono) The injections $\inl\colon A\to A+B$ and $\inr\colon B\to A+B$ of binary coproducts in $\C$ are monic.

\medskip\noindent (\hpresint) The functor $H$ preserves binary intersections (i.e.\ pullbacks of monomorphisms).

\medskip\noindent We will tag all results below with the conditions they require (on top of the global conditions given by \autoref{asm}). For example, \autoref{prop:wf-coalgebra-coproduct} requires both (\injmono) and (\hpresint).

\begin{rem}\label{rem:conditions}
The two extra conditions (\injmono) and (\hpresint) are fairly mild:
\begin{enumerate}
\item Condition (\injmono) holds in every extensive category with finite limits~\cite[Prop.~3.3]{cw93}. Recall that a category $\C$ is \emph{extensive}~\cite{cw93} if it has finite coproducts and for any pair of objects~$A$,~$B$ the functor $\C/A\times \C/B \xto{+} \C/(A+B)$ on slice categories sending $(f\colon X\to A)$, $(g\colon Y\to B)$ to $(f+g\colon X+Y\to A+B)$ is an equivalence of categories. Informally, this means that coproducts behave like disjoint unions. The categories of sets, posets, and nominal sets are extensive, as is every elementary topos~\cite{Johnstone}. There are also many non-extensive categories with (\injmono), including the categories of vector spaces or convex sets (\autoref{sec:convex}).
\item\label{rem:conditions:set-func-pres-int} Condition (\hpresint) `almost' holds for every functor $H\colon \Set\to \Set$. By a classical result due to Trnkov\'a~\cite{trnkova71} (see also~\cite[Prop.~8.1.13]{amm25}), every set functor~$H$ can be modified, by changing its action on the empty set and functions with empty domain, to a functor $\widetilde{H}$ that preserves binary intersections. Since $\Coalg(H)\cong \Coalg(\widetilde{H})$, one can thus assume that $H$ preserves binary intersections without loss of generality for coalgebraic results.
\end{enumerate} 
\end{rem}

\begin{rem}\label{rem:int-implies-mono}
\begin{enumerate}
\item Condition ($\injmono$) implies that the initial object $0$ is simple: given $X\in \C$, the unique morphism $0\to X$ is the coproduct injection $\inl\colon 0\to 0+X\cong X$, hence monic. 
\item Condition (\hpresint) implies that $H$ preserves monomorphisms. This follows from the fact that $m\colon A\to B$ is monic iff the pair $\id_A,\id_A$ is a pullback of $m,m$.
\end{enumerate}

\end{rem}

We start with a simple construction, the \emph{coproduct extension} of a coalgebra. It appeared earlier in the theory of iterative algebras~\cite{amv06,urbat17} and underlies both our proof of the coalgebraic \Konig's lemma and the initial algebra constructions presented in \autoref{sec:initial-algebras}. 

\begin{construction}[Coproduct Extension]\label{cons:cp}
The \emph{coproduct extension} of an $H$-coalgebra $(C,c)$ by a morphism $p\colon X\to HC$ is the coalgebra $(C+X,c_p)$ whose structure is given by
\[ c_p \, \equiv \, (\,C+X \xto{[c,p]} HC \xto{H\inl} H(C+X)\,).   \]
\end{construction}
Let us immediately note an expected property of coproduct extensions:

\begin{lemma}\label{lem:coproduct-inl-coalg-mor} The coproduct injection $\inl\colon (C,c)\to (C+X,c_p)$ is a coalgebra morphism.
\end{lemma}

\begin{proofappendix}{lem:coproduct-inl-coalg-mor}
This is shown by the commutative diagram below:
\[
\begin{tikzcd}[baseline=(HC.base)]
C \ar{dr}{c} \ar{r}{\inl} \ar{dd}[swap]{c} & C+X \ar[shiftarr={xshift=15mm}]{dd}{c_p} \ar{d}{[c,p]} \\
& HC \ar{d}{H\inl} \\
|[alias=HC]|
HC \ar{r}{H\inl} & H(C+X)
\end{tikzcd}
\tag*{\qedhere}
\]
\end{proofappendix}

Informally, a coproduct extension adds to a given coalgebra $(C,c)$ a set $X$ of new states and lets every new state transition arbitrarily into states from $C$, while leaving transitions of states in $C$ unchanged. This construction should not introduce any infinite paths to the given coalgebra, which is confirmed by the following result: 

\begin{proposition}[\injmono, \hpresint]\label{prop:wf-coalgebra-coproduct}
For every well-founded coalgebra $(C,c)$ and every morphism $p\colon X\to HC$, the coproduct extension $(C+X,c_p)$  is a well-founded coalgebra.
\end{proposition}

\begin{proof}[Proof sketch]
Let $m\colon (S,s)\monoto (C+X,c_p)$ be a cartesian subcoalgebra. We need to prove that $m$ is an isomorphism. Form the two intersections in $\C$ shown below; note that the coproduct injections $\inl$ and $\inr$ are monic by assumption (\injmono).
\begin{equation*}
\begin{tikzcd}
S\cap C \pullbackangle{-45} \ar[tail]{d}[swap]{l_{S\cap C}} \ar[tail]{r}{r_{S\cap C}} & C \ar[tail]{d}{\inl} \\
S \ar[tail]{r}{m} & C+X
\end{tikzcd}
\qquad
\begin{tikzcd}
S\cap X \pullbackangle{-45} \ar[tail]{d}[swap]{l_{S\cap X}} \ar[tail]{r}{r_{S\cap X}} & X \ar[tail]{d}{\inr} \\
S \ar[tail]{r}{m} & C+X
\end{tikzcd}
\end{equation*}
One can show that both $r_{S\cap C}$ and  $r_{S\cap X}$ are isomorphisms. In the case of $r_{S\cap C}$, this follows from the observation that the subobject $r_{S\cap C}\colon S\cap C \monoto C$ carries a cartesian subcoalgebra of the well-founded coalgebra $(C,c)$. We conclude that $m$ is a split epimorphism:
\[ m\circ [l_{S\cap C} \circ r_{S\cap C}^{-1}, l_{S\cap X}\circ r_{S\cap X}^{-1}] = [\inl,\inr]=\id_{C+X}, \]
and since $m$ is also monic, it is an isomorphism as required. 
\end{proof}

\begin{proofappendix}{prop:wf-coalgebra-coproduct}
Let $m\colon (S,s)\monoto (C+X,c_p)$ be a cartesian subcoalgebra. We need to prove that $m$ is an isomorphism. Form the following intersections in $\C$ (the coproduct injections $\inl$ and $\inr$ are monic by assumption \injmono):
\begin{equation}\label{eq:sc-int}
\begin{tikzcd}
S\cap C \pullbackangle{-45} \ar[tail]{d}[swap]{l_{S\cap C}} \ar[tail]{r}{r_{S\cap C}} & C \ar[tail]{d}{\inl} \\
S \ar[tail]{r}{m} & C+X
\end{tikzcd}
\end{equation}
\begin{equation}\label{eq:sx-int}
\begin{tikzcd}
S\cap X \pullbackangle{-45} \ar[tail]{d}[swap]{l_{S\cap X}} \ar[tail]{r}{r_{S\cap X}} & X \ar[tail]{d}{\inr} \\
S \ar[tail]{r}{m} & C+X
\end{tikzcd}
\end{equation}
Below we shall prove the following:
\begin{enumerate}
\item\label{s1} $r_{S\cap C}$ is an isomorphism.
\item\label{s2} $r_{S\cap X}$ is an isomorphism.
\end{enumerate}
These statements imply that $m$ is a split epimorphism:
\[ m\circ [l_{S\cap C} \circ r_{S\cap C}^{-1}, l_{S\cap X}\circ r_{S\cap X}^{-1}] = [\inl,\inr]=\id_{C+X}, \]
and since $m$ is also monic, it is an isomorphism as required. 

It thus only remains to prove the above two claims.

\medskip\noindent \emph{Proof of \ref{s1}}: The following diagram commutes:
\begin{equation}\label{eq:comm-diag}
\begin{tikzcd}
S\cap C \ar[tail]{dd}[swap]{r_{S\cap C}} \ar[tail]{rr}{l_{S\cap C}}  && S \ar{d}{s} \ar[tail]{dl}[swap]{m} \\
& C+X \ar{d}{[c,p]} \ar{dr}{c_p} & HS \ar[tail]{d}{Hm} \\
C \ar{ur}{\inl} \ar{r}{c} & HC \ar{r}{H\inl} & H(C+X)
\end{tikzcd}
\end{equation}
Therefore, by \eqref{eq:sc-int} and since $H$ preserves intersections, there exists a unique morphism $s_C$ making the left-hand and upper part of the diagram below commute:
\begin{equation}\label{eq:def-ol-sc}
\begin{tikzcd}
S\cap C \ar[tail]{dd}[swap]{r_{S\cap C}} \ar[tail]{rr}{l_{S\cap C}}  \ar[dashed]{dr}{s_C} && S \ar{d}{s} \\
& H(S\cap C)  \pullbackangle{-45}\ar[tail]{r}{Hl_{S\cap C}} \ar[tail]{d}{Hr_{S\cap C}} & HS \ar[tail]{d}{Hm} \\
C \ar{r}{c} & HC \ar{r}{H\inl} & H(C+X)
\end{tikzcd}
\end{equation}
We claim that $r_{S\cap C}\colon (S\cap C, s_C) \monoto (C,c)$ is a cartesian subcoalgebra, which means that the following diagram is a pullback:
\begin{equation}\label{eq:sc-pullback}
\begin{tikzcd}
S \cap C \ar[tail]{r}{s_C} \ar[tail]{d}[swap]{r_{S\cap C}} \pullbackangle{-45} & H(S\cap C) \ar[tail]{d}{Hr_{S\cap C}}\\
C \ar{r}{c} & HC
\end{tikzcd}
\end{equation}
Thus let $f\colon Z\to H(S\cap C)$ and $g\colon Z\to C$ be morphisms with
\begin{equation}\label{eq:def-fg}
Hr_{S\cap C}\circ f = c\circ g.
\end{equation}
Then the diagram below commutes:
\begin{equation}\label{eq:comm-diag2}
\begin{tikzcd}
Z \ar{rrr}{f} \ar{dd}[swap]{g} & & & H(S\cap C) \ar[tail]{d}{Hl_{S\cap C}} \ar[tail]{ddl}[swap]{Hr_{S\cap C}}  \\
& & & HS \ar[tail]{d}{Hm} \\
  C \ar[rounded corners,to path={
    ([xshift=2mm]\tikztostart.north)
    -- ++(0,4mm)
    -- ([yshift=4mm,xshift=-1mm]\tikztotarget.north) \tikztonodes
    -- ([xshift=-1mm]\tikztotarget.north)
  }]{rr}{c} \ar{r}{\inl} & C+X \ar{r}{[c,p]} & HC \ar{r}{H\inl} & H(C+X)
\end{tikzcd}
\end{equation}
Therefore, there exists a unique morphism $j$ making the left-hand and upper part of the following diagram commute:
\begin{equation}\label{eq:def-j}
\begin{tikzcd}
Z \ar[dashed]{dr}{j} \ar{rrr}{f} \ar{dd}[swap]{g} & & & H(S\cap C) \ar[tail]{d}{Hl_{S\cap C}}  \\
& S \pullbackangle{-45} \ar[tail]{d}[swap]{m} \ar{rr}{s} & & HS \ar[tail]{d}{Hm} \\
C \ar{r}{\inl} & C+X \ar{r}{[c,p]} & HC \ar{r}{H\inl} & H(C+X)
\end{tikzcd}
\end{equation}
This in turn yields a unique morphism $k$ making the left-hand and upper part of the diagram
\begin{equation}\label{eq:def-k}
\begin{tikzcd}
Z \ar[dashed]{dr}{k} \ar[bend left=20]{drr}{g} \ar[bend right=20]{ddr}[swap]{j} & &  \\
& S\cap C \pullbackangle{-45} \ar[tail]{r}{r_{S\cap C}} \ar[tail]{d}[swap]{l_{S\cap C}} & C \ar[tail]{d}{\inl} \\
& S \ar[tail]{r}{m} & C+X
\end{tikzcd}
\end{equation}
commute. Then $k$ also makes the left-hand and upper part of the diagram below commute:
\begin{equation}\label{eq:k-pullback}
\begin{tikzcd}
Z \ar[bend left=20]{drr}{f} \ar[bend right=20]{ddr}[swap]{g} \ar{dr}{k} && \\
& S \cap C \ar[tail]{r}{s_C} \ar[tail]{d}[swap]{r_{S\cap C}} & H(S\cap C) \ar[tail]{d}{Hr_{S\cap C}}\\
& C \ar{r}{c} & HC
\end{tikzcd}
\end{equation}
Indeed, the left-hand part commutes by definition of $k$, and the upper part commutes because it does so when post-composed with the monomorphism $Hr_{S\cap C}$. Moreover, since $r_{S\cap C}$ is monic, $k$ is the unique morphism making both parts commute. This proves that \eqref{eq:sc-pullback} is a pullback, as claimed. Since $(C,c)$ is well-founded, we conclude that $r_{S\cap C}$ is an isomorphism.

\medskip\noindent \emph{Proof of \ref{s2}}:
The following diagram commutes:
\begin{equation}\label{eq:comm-diag3}
\begin{tikzcd}[column sep=30]
X \ar{rr}{p} \ar{ddd}[swap]{\inr}  && HC \ar[bend right=10, equals]{dddl} \ar{d}{Hr_{S\cap C}^{-1}} \\
&& H(S\cap C) \ar{d}{Hl_{S\cap C}} \\
&& HS \ar{d}{Hm}  \\
C+X \ar{r}{[c,p]} & HC \ar{r}{H\inl} & H(C+X)
\end{tikzcd}
\end{equation}
Therefore, there exists a unique morphism $n$ making the left-hand and upper part of the diagram below commute:
\begin{equation}\label{eq:comm-diag4}
\begin{tikzcd}[column sep=30]
X \ar[dashed]{ddr}{n} \ar{rrr}{p} \ar[bend right=10]{dddr}[swap]{\inr}  &&& HC \ar{d}{Hr_{S\cap C}^{-1}} \\
&&& H(S\cap C) \ar{d}{Hl_{S\cap C}} \\
& S \ar{rr}{s} \ar[tail]{d}[swap]{m} \pullbackangle{-45} && HS \ar{d}{Hm}  \\
&C+X \ar{r}{[c,p]} & HC \ar{r}{H\inl} & H(C+X)
\end{tikzcd}
\end{equation}
This in turn yields a unique morphism $o$ making the left-hand and upper part of the following diagram commute:
\begin{equation}
\begin{tikzcd}
X \ar[shiftarr={yshift=15}]{rr}{n} \ar[dashed]{r}{o} \ar{dr}[swap]{\id} & S\cap X \pullbackangle{-45} \ar[tail]{r}{l_{S\cap X}} \ar[tail]{d}[swap]{r_{S\cap X}} & S \ar[tail]{d}{m} \\
& X \ar{r}{\inr} & C+X 
\end{tikzcd}
\end{equation}
Thus the monomorphism $r_{S\cap X}$ is a split epimorphism (with splitting $o$), and so it is an isomorphism.
\end{proofappendix}\begin{figure}  \begin{minipage}[b]{.40\textwidth}\begin{tikzcd}[column sep=4mm,every cell/.append style={inner xsep=1pt}]
(C_i,c_i)
  \ar[shiftarr={yshift=8mm}, tail]{rr}{m_i} \ar[tail]{r}{n_i}
& (\bigvee_{i\in I} C_i, \ol{c}) \ar[tail]{r}{m}
& (C,c)  
\end{tikzcd}
  \caption{Join of subcoalgebras.}
  \label{fig:joinSubcoalgebras}
  \end{minipage}  \begin{minipage}[b]{.30\textwidth}\begin{tikzcd}[outer sep=0pt,row sep=6mm]
\bigvee C_i \ar{r}{\overline{c}} \pullbackangle{-45} \ar[tail]{d}[swap]{m} & H(\bigvee C_i) \ar[tail]{d}{Hm} \\
C \ar{r}{c} & HC
\end{tikzcd}
  \caption{Cartesianess.}
  \label{fig:joinCartesian}
  \end{minipage}  \begin{minipage}[b]{.30\textwidth}
    \centering
  \begin{tikzcd}[row sep=3mm,column sep=4mm]
      \phantom{-}1 \arrow{r} \arrow{dr}
      & \phantom{-}2 \arrow{r} \arrow{dr}
      & \phantom{-}3 \arrow{r} \arrow{dr}
      & \cdots
      \\
      -1 \arrow{r} \arrow{ur}
      & -2 \arrow{r} \arrow{ur}
      & -3 \arrow{r} \arrow{ur}
      & \cdots
    \end{tikzcd}
    \\[2mm]
    $k \mapsto (-|k|{-}1, |k|{+}1)$
  \caption{A coalgebra}
  \label{fig:coalgExample}
  \end{minipage}\end{figure}

Recall the classical \Konig's lemma: In a finitely branching well-founded graph (equivalently, a well-founded $\Powf$-coalgebra), every node lies in some finite subcoalgebra. The following theorem generalizes \Konig's lemma to coalgebras over \lfp categories:


\begin{theorem}[Coalgebraic \Konig's Lemma, \injmono, \hpresint]\label{thm:koenig}
Every well-founded $H$-coalgebra is the directed join of its fg-carried well-founded subcoalgebras.
\end{theorem}

Here the join is formed in the complete lattice of subcoalgebras of the given coalgebra. As a minor subtlety, we note that subcoalgebras of well-founded coalgebras generally need not be well-founded; this would require stronger assumptions~\cite[Prop.~8.4.7]{amm25}. Therefore, it is crucial that \autoref{thm:koenig} considers \emph{well-founded} (not arbitrary) subcoalgebras. 


\begin{proof}[Proof sketch]
Let $(C,c)$ be a well-founded coalgebra, and let $m_i\colon (C_i,c_i)\monoto (C,c)$ ($i\in I$) be the family of all fg-carried well-founded subcoalgebras of $(C,c)$. Under the usual subobject ordering, the family forms a directed poset; this follows easily from the closure properties of $\C_\fg$ (\autoref{lem:fp-fp-fin-colimits}) and $\Coalg_{\wf}(H)$ (\iref{prop:wf}{closed-colim-quotients}). Form the (directed) join of all $m_i$ in
in the complete lattice of subcoalgebras of $(C,c)$ (see \autoref{fig:joinSubcoalgebras}). This join is computed at the level of the subobject lattice of $C\in \C$ (\autoref{lem:colimits-coalg}). Since directed joins 


\noindent of subobjects in \lfp categories are directed colimits (\autoref{lem:lfp-directed-unions}), we see that $(n_i)_{i\in I}$ is a colimit cocone. It suffices to show that $m$ is a cartesian subcoalgebra, since this implies that $m$ is an isomorphism and hence proves the theorem. 

For the proof that $m$ is cartesian (\autoref{fig:joinCartesian}), one needs to show that the commutative square below forms a pullback, that is, given $f\colon X\to H(\bigvee C_i)$ and $g\colon X\to C$ with
$Hm\circ f = c\circ g$, one
needs to construct the mediating morphism $k\colon X\to \bigvee C_i$. This is achieved by first constructing $k$ for the case $X\in \C_\fg$. Here one factorizes the morphism $f$ as $X\xto{p} HC_i \xto{Hn_i} H(\bigvee_i C_i)$ for some~$p$ and $i$, using that $(Hn_i)$ is a colimit cocone since the functor $H$ is finitary, and then analyses the coproduct extension $(C_i+X,(c_i)_p)$. The general case $X\in \C$ follows from the case 
of finitely generated objects by expressing $X$ as the colimit of its finitely generated subobjects. This is the point in the proof where we need that $\C$ is \lfp. 
\end{proof}

\begin{proofappendix}{thm:koenig}
Let $(C,c)$ be a well-founded coalgebra and let $\Sub_{\fg,\wf}(C,c)$ denote the poset of all its fg-carried well-founded subcoalgebras. Note that the poset $\Sub_{\fg,\wf}(C,c)$ is small because $\C$ is well-powered (\autoref{lem:lfp-complete-wp}). Moreover, it is directed, in fact a lattice: Given a finite family $m_i\colon (C_i,c_i)\monoto (C,c)$ ($i\in I_0$) of fg-carried well-founded subcoalgebras, its join $m\colon \bigvee_{i\in I_0} (C_i,c_i)\monoto (C,c)$ in $\Sub_{\fg,\wf}(C,c)$  is the image of $[m_i]_{i\in I_0}\colon \coprod_{i\in I_0}\colon (C_i,c_i)\to (C,c)$. Note that the coalgebra $\bigvee_{i\in I_0} (C_i,c_i)$ actually lies in $\Coalg_{\fg,\wf}(H)\hookrightarrow \Coalg(H)$ because the latter subcategory is closed under finite coproducts (\autoref{lem:fp-fp-fin-colimits}) and strong quotients (\autoref{prop:wf}). 

Now let $m_i\colon (C_i,c_i)\monoto (C,c)$ ($i\in I$) be the family of all fg-carried well-founded subcoalgebras of $(C,c)$. Form its join in the complete lattice of all subcoalgebras of $(C,c)$:
\[
\begin{tikzcd}
(C_i,c_i) \ar[shiftarr={yshift=6mm}, tail]{rr}{m_i} \ar[tail]{r}{n_i} & (\bigvee_{i\in I} C_i, \ol{c}) \ar[tail]{r}{m} & (C,c)  
\end{tikzcd}
\]
This join is directed, as argued above, and formed at the level of $\C$ (\autoref{lem:colimits-coalg}). Since directed joins of subobjects in \lfp categories are directed colimits (\autoref{lem:lfp-directed-unions}), we see that $(n_i)_{i\in I}$ is a colimit cocone. We will show that $m$ is a cartesian subcoalgebra, which implies that $m$ is an isomorphism and hence proves the theorem.

To prove that $m$ is cartesian, we need to show that the following square is a pullback:
\begin{equation}\label{eq:pullback}
\begin{tikzcd}
\bigvee C_i \ar{r}{\overline{c}} \pullbackangle{-45} \ar[tail]{d}[swap]{m} & H(\bigvee C_i) \ar[tail]{d}{Hm} \\
C \ar{r}{c} & HC
\end{tikzcd}
\end{equation}

Thus suppose that $f\colon X\to H(\bigvee C_i)$ and $g\colon X\to C$ are morphisms with
\[ Hm\circ f = c\circ g.  \]
Our task is to construct the mediating morphism $k\colon X\to \bigvee C_i$. 

\begin{enumerate}
\item Assume first that $X\in \C_\fg$. Since $H$ is finitary and preserves monos, $(Hn_i)$ is a colimit cocone of monomorphisms and so $f$ factorizes as
\[ 
\begin{tikzcd}
X \ar{r}{f} \ar{dr}[swap]{p} & H(\bigvee C_i) \\
& HC_i \ar[tail]{u}[swap]{Hn_i}
\end{tikzcd}
 \]
for some $i\in I$ and $p\colon X\to HC_i$. Form the coproduct extension $(C_i+X,(c_i)_p)$ of the coalgebra $(C_i,c_i)$ by $p$ (\autoref{cons:cp}). Then
\[ [m_i,g]\colon (C_i+X,(c_i)_p)\to (C,c)\]
is a coalgebra morphism, as the diagram below commutes:
\[
\begin{tikzcd}
C_i+X \ar{rr}{(c_i)_p} \ar{dr}{[c_i,p]} \ar{dd}[swap]{[m_i,g]} && H(C_i+X) \ar{dd}{H[m_i,g]}  \\
& HC_i \ar{ur}{H\inl} \ar{dr}{Hm_i}  & \\
C \ar{rr}{c} && HC 
\end{tikzcd}
\]
For commutativity of the second component of the lower left part, we compute
\[ Hm_i\circ p = Hm\circ Hn_i \circ p = Hm\circ f = c\circ g. \]
Form the image factorization of $[m_i,g]$. Since the coalgebra $(C_i+X, (c_i)_p)$ is well-founded (\autoref{prop:wf-coalgebra-coproduct}) and $\Coalg_{\fg,\wf}(H)$ closed under strong quotients (\autoref{lem:fp-fp-fin-colimits}, \iref{prop:wf}{closed-colim-quotients}), the image of $[m_i,g]$ is an fg-carried well-founded subcoalgebra of $(C,c)$ and thus of the form $m_j\colon (C_j,c_j)\monoto (C,c)$ for some $j\in J$:
\[  
\begin{tikzcd}
(C_i+X,(c_i)_p) \ar[shiftarr={yshift=15}]{rr}{[m_i,g]} \ar[two heads]{r}{e} & (C_j,c_j) \ar[tail]{r}{m_j} & (C,c)
\end{tikzcd}
\]
We now define $k\colon X\to \bigvee_i C_i$ to be the composite
\[ k = (\, X \xto{\inr} C_i+X \xto{e} C_j \xto{n_j} \bigvee_i C_i  \,). \]
Then $k$ is the desired mediating morphism, i.e. both the left-hand and upper part of the diagram below commute:
\begin{equation}\label{eq:pb-proof}
\begin{tikzcd}
X \ar[bend left=20]{drr}{f} \ar[bend right=20]{ddr}[swap]{g} \ar{dr}{k}  && \\
& \bigvee C_i \ar{r}{\overline{c}} \pullbackangle{-45} \ar[tail]{d}[swap]{m} & H(\bigvee C_i) \ar[tail]{d}{Hm} \\
& C \ar{r}{c} & HC
\end{tikzcd}
\end{equation}
Indeed, commutativity of the left-hand part is witnessed by the following diagram:
\[
\begin{tikzcd}
X  \ar{d}[swap]{g} \ar[shiftarr={yshift=7mm}]{rrr}{k} \ar{r}{\inr} & C_i+X \ar{d}{[m_i,g]} \ar{r}{e} & C_j \ar{d}{m_j} \ar{r}{n_i} & \bigvee_i C_i \ar{d}{m} \\
C \ar[equals]{r} & C \ar[equals]{r} & C \ar[equals]{r} & C 
\end{tikzcd}
\]
Then also the upper part of \eqref{eq:pb-proof} commutes because it does so when postcomposed with the monomorphism $Hm$. The uniqueness of $k$ follows from $m$ being monic.
\item Now let $X\in \C$ be an arbitrary object. Since $\C$ is \lfp, we can express $X$ as the directed join of the family $x_j\colon X_j\monoto X$ ($j\in J$) of its finitely generated subobjects. By part 1 above, for each $j\in J$ there exists a morphism $k_j\colon X_j\to \bigvee_i C_i$ such that
\[ m\circ k_j = g\circ x_j. \]
The morphisms $k_j$ form a cocone over the diagram of $X_j$'s, so there exists a unique $k\colon X\to \bigvee_i C_i$ such that $k\circ x_j = k_j$ for all $j\in J$. This implies
\[ m\circ k\circ x_j = m\circ k_j = g\circ x_j  \]
and so
\[ m\circ k = g \]
because the colimit injections $x_j$ are jointly epic. Thus $k$ is the desired mediating morphism witnessing that \eqref{eq:pullback} is a pullback.\qedhere
\end{enumerate}
\end{proofappendix}


As an immediate consequence, we get a property of the category of well-founded coalgebras:

\begin{corollary}[\smash{\injmono}, \hpresint]\label{cor:wfIsLFP}
If $\C_\fp=\C_\fg$, then $\Coalg_\wf(H)$ is \lfp. 
\end{corollary}

\begin{proof}
The category $\Coalg_\wf(H)$ is cocomplete because it is closed under colimits in $\Coalg(H)$ (\iref{prop:wf}{closed-colim-quotients}). All coalgebras in $\Coalg_{\fp,\wf}(H)$ 
are finitely presentable objects of $\Coalg_\wf(H)$; this follows from the fact that all coalgebras in $\Coalg_\fp(H)$ are finitely presentable objects of $\Coalg(H)$ (\autoref{lem:fp-fg-coalg}) and that $\Coalg_\wf(H)$ is closed under (filtered) colimits in $\Coalg(H)$. By \autoref{thm:koenig}, every object of $\Coalg_\wf(H)$ is a filtered colimit of objects in $\Coalg_{\fp,\wf}(H)=\Coalg_{\fg,\wf}(H)$. Thus $\Coalg_\wf(H)$ is \lfp. 
\end{proof}

Without the restrictive assumption that $\C_\fp=\C_\fg$, we get a slightly weaker result. It makes use of the following characterization of locally presentable categories (cf.\ \autoref{rem:loc-pres}):

\begin{theorem}[{Local Generation Theorem~\cite[Thm.~1.70]{adamek_rosicky_1994}}] A category is locally presentable iff it is cocomplete and, for some regular cardinal $\lambda$, there exists a set $\mathcal{A}$ of $\lambda$-generated objects such that every object is the $\lambda$-directed colimit of the diagram of its subobjects from $\mathcal{A}$. 
\end{theorem}

Note that the `if' direction of the theorem does not claim that the category is locally $\lambda$-presentable. We now instantiate this theorem to the category $\Coalg_\wf(H)$, $\lambda=\omega$, and $\mathcal{A}=\Coalg_{\fg,\wf}(H)$. The coalgebras in $\A$ are finitely generated objects of $\Coalg_\wf(H)$ because they are finitely generated in $\Coalg(H)$ (\autoref{lem:fp-fg-coalg}) and  $\Coalg_\wf(H)$ is closed under filtered colimits in $\Coalg(H)$ (\iref{prop:wf}{closed-colim-quotients}). Every coalgebra in $\Coalg_\wf(H)$ is the directed colimit of its subobjects from $\A$ (\autoref{thm:koenig}). Thus, the Local Generation Theorem yields:

\begin{corollary}[\smash{\injmono}, \hpresint]
The category $\Coalg_\wf(H)$ is locally presentable.
\end{corollary}
We conclude this section with the observation that the coalgebraic \Konig's lemma no longer holds when well-founded coalgebras are replaced with recursive ones, even in $\Set$:

\begin{example}[\Konig's Lemma Fails for Recursive Coalgebras]\label{ex:konig-lemma-recursive}
  Consider the set functor
  \[
    HX = \set{*} + \set{\,(x_1,x_2)\in X\times X \mid x_1\neq x_2\,}
  \]
  which is the quotient of $X\times X$ identifying all elements on the diagonal with a constant~$*$. Clearly, $H$ is finitary and preserves binary intersections. Let $(C,c)$ be the coalgebra with state space $C := \Z\setminus\set{0}$ and the transition structure
  $c\colon C \to HC$, $c(k) = (-|k|-1, |k|+1)$ visualized in \autoref{fig:coalgExample}.
                                    The only finite subcoalgebra of $(C,c)$ is the empty coalgebra; thus, $(C,c)$ is not the union of its finite recursive subcoalgebras. However, $(C,c)$ is recursive: For every $H$-algebra $(A,a)$, the constant map $h\colon (C,c)\to (A,a)$ defined by $h(k) :=
  a(*)$ is the unique coalgebra-to-algebra morphism. Indeed, $h$ is such a morphism because $(Hh\circ  c)(k) = *\in HA$ and so $(a\circ Hh\circ c)(k)=a(*)=h(k)$ for all $k\in C$. To prove uniqueness, suppose that $g\colon (C,c)\to (A,a)$ is a coalgebra-to-algebra morphism, and let $k\in C$. Since $c(|k|+1) = c(-|k|-1)$,
  we necessarily have $g(|k|+1) = g(-|k|-1)$ by $g = a\circ HG\circ c$. Thus, $Hg(c(k)) = * \in HA$ and again by $g = a\circ HG\circ c$, we obtain
\[ g(k) = (a\circ Hg\circ c)(k) = a(*) = h(k).\]
\twnote[inline]{}
\end{example}
\begin{proofappendix}[Details for]{ex:konig-lemma-recursive}
  For the detailed verification of the uniqueness of $h$, consider some coalgebra-to-algebra morphism $g\colon (C,c)\to (A,a)$ and $k\in C$.
  We have $c(|k|+1) = c(-|k|-1)$ and so
    \[
      g(|k|+1)
      = (a \circ Hg \circ c)(|k|+1)
      = (a \circ Hg\circ c)(-|k|-1)
      = g(-|k|-1).
    \]
    The definition of $c$ thus implies
    $Hg(c(k)) = * \in HA$, which proves
  \[ g(k) = (a\circ Hg\circ c)(k) = a(*) = h(k).
  \tag*{\qedhere\qed}
  \] 
\end{proofappendix}

\section{Applications of the Coalgebraic \Konig's Lemma}\label{sec:applications}
We present a few selected instantiations of our coalgebraic \Konig's lemma (\autoref{thm:koenig}).

\subsection{Coalgebras for Set Functors}\label{sec:finitary-set}
For the category $\C=\Set$, we recover a recent result by Ad\'amek et al.~\cite[Prop,~9.2.19]{amm25} as a special case of \autoref{thm:koenig}:

\begin{theorem}[\Konig's Lemma for Coalgebras in $\Set$]\label{thm:koenig-coalgebra-set} If $H$ is a finitary set functor, then every state of a well-founded $H$-coalgebra lies in some finite subcoalgebra.
\end{theorem}

Indeed, by \iref{rem:conditions}{set-func-pres-int} we can safely assume that a finitary set functor $H$ preserves binary intersections (by modifying $H$ on the empty set). Thus \autoref{thm:koenig} applies.

\subsection{Graphs in a Topos}\label{sec:topos}
While \autoref{thm:koenig-coalgebra-set} extends \Konig's lemma from graphs to coalgebras in $\Set$, we now take an orthogonal direction: we stick with graphs but interpret them in a \lfp \emph{elementary topos}~$\C$~\cite{Johnstone_short}. (We assume basic familiarity with topos theory in this subsection.) Examples of \lfp toposes include the topos of sets, nominal sets (see \autoref{sec:nom}), any presheaf topos $\Set^\D$, or any coherent Grothendieck topos~\cite[Rem.~D.3.3.12]{Johnstone}.  

Graphs in a topos are coalgebras for the covariant power object functor $\Pow\colon \C\to \C$~\cite[Sec.~A.2.3]{Johnstone}; the latter maps each object $C$ to its power object $\Pow C=\Omega^C$, where $\Omega$ is the subobject classifier. By cartesian closure, a coalgebra $c\colon C\to \Pow C=\Omega^C$ corresponds to a morphism from $C\times C $ to $\Omega$ and thus to a subobject of $C\times C$, that is, a \emph{relation} on $C$. The role of the finite power set functor (modelling finitely branching graphs) is now played by the finitary coreflection $\Powf$ of $\Pow$ (\autoref{lem:finitary-coreflection}). Note that $\Powf$ is generally different from the \emph{Kuratowski functor} $\mathcal{K}\colon \C\to \C$, which captures Kuratowski-finite subobjects~\cite[Thm.~9.16]{Johnstone_short}. 

To see that $\Powf$-coalgebras indeed model `finite' branching, let $i\colon \Powf\to \Pow$ denote the co-universal natural transformation witnessing that $\Powf$ is the finitary coreflection of $\Pow$. Every $\Powf$-coalgebra $c\colon C\to \Powf C$ can be viewed as a $\Pow$-coalgebra $i_C\circ c\colon C\to \Pow C$ and therefore as a
 relation $m_c\colon E_c\monoto C\times C$. Given $s\colon I\to C$ we let $m_{c,I}\colon E_{c,I} \monoto I\times C$ denote the preimage of $m_c$ under $s\times \id$:
\[
\begin{tikzcd}[row sep=4mm,outer sep=0pt]
E\smash{{}_{c,I}} \pullbackangle{-45} \ar{r} \ar[tail]{d}[swap]{m_{c,I}} & E\smash{{}_c} \ar[tail]{d}{m_c} \\
I\times C \ar{r}{s\times \id} & C\times C
\end{tikzcd}
\]
 We then have the following simple result. Informally, it expresses that if we consider a finitely indexed family $(s_i)_{i\in I}$ of states in $C$, then their successors can also be finitely indexed. 


\begin{lemma}\label{lem:powf-coalg-topos}
Let $\C$ be a \lfp topos, and let $c\colon C\to \Powf C$ be a coalgebra for $\Powf$. Then for every $s\colon I\to C$ ($I\in \C_\fp$) there exists $t\colon J \to C$ ($J\in \C_\fp$) such that $m_{c,I}\colon E_{c,s}\monoto I\times C$ factorizes through $\id\times t\colon I\times J\to I\times C$. 
\end{lemma}

\begin{proofappendix}{lem:powf-coalg-topos}
 Recall that $\Powf C$ is the colimit of the canonical diagram $D_C\colon\C_\fp/C\to \C$ mapping $(t\colon J\to C)$ to $\Pow J$ (\autoref{lem:finitary-coreflection}). We denote the colimit injection for $t$ by $e_t\colon \Pow J \to \Powf C$.

Now let $s\colon I\to C$ where $I\in \C_\fp$. Since $(e_t)$ is a filtered colimit, the morphism $c\circ s$ factorizes through $e_t$ for some $t\colon J\to C$ ($J\in \C_\fp$), that is, there exists a morphism $I\to \Pow J$ making the left-hand part of the diagram below commute.
\[
\begin{tikzcd}
&& \Pow J \ar{dr}{\Pow t} \ar{d}{e_t} & \\
I \ar[bend left=20]{urr}{\exists} \ar{r}{s} & C \ar{r}{c} & \Powf C \ar{r}{i_C} & \Pow C
\end{tikzcd}
\]
Since the right-hand triangle also commutes, we see that the morphism $i_C\circ c \circ s$ factorizes through $\Pow t$. This morphism corresponds to the subobject $m_{c,I}\colon E_{c,I}\monoto S\times C$, so the statement of the lemma follows.
\end{proofappendix}

Applying the coalgebraic \Konig's lemma to the case of $\Powf$-coalgebras yields:

\begin{theorem}[\Konig's Lemma for Graphs in a Topos]\label{thm:koenigTopos}
  Let $\C$ be a \lfp topos. Every well-founded $\Powf$-coalgebra is the join of its fg-carried well-founded subcoalgebras.  
\end{theorem}

\begin{proofappendix}{thm:koenigTopos}
We show that the category $\C$ and the functor $\Powf$ satisfy the conditions of \autoref{thm:koenig}. Every topos is an extensive category, so (\injmono) holds (\autoref{rem:conditions}). The functor $\Powf$ is finitary by definition. Moreover, the power object functor $\Pow$ preserves binary intersections~\cite[Thm.~5.9]{osius74}, and this property carries over to its finitary coreflection $\Powf$ because the latter is objectwise constructed via canonical filtered colimits (\autoref{lem:finitary-coreflection}), and filtered colimits commute with finite limits (in particular intersections) in \lfp categories~\cite[Prop.~1.59]{adamek_rosicky_1994}.
\end{proofappendix}

\subsection{Binary Trees}
An important special case of \Konig's lemma, studied in constructive mathematics~\cite{di21}, is its restriction to binary trees -- the \emph{weak \Konig's lemma}. Here a {binary} tree is one where every node has at most two children, and children are ordered, that is, binary trees form coalgebras for $HX = X\times X + X + 1$ on $\Set$. For \lfp categories with nice coproducts (cf.~\autoref{rem:conditions}), the generalized \Konig's lemma also has a weak version:

\begin{theorem}[Coalgebraic Weak \Konig's Lemma]\label{cor:weak-koenigs-lemma}
Let $\C$ be a \lfp and extensive category. Then every well-founded coalgebra for the functor $HX = X\times X + X + 1$ on $\C$ is the join of its fg-carried well-founded subcoalgebras.  
\end{theorem}
This theorem extends to polynomial functors $HX = \coprod_{i\in I} X^{n_i}$ where $I$ is finite and $n_i\in \Nat$.

\begin{proofappendix}{cor:weak-koenigs-lemma}
Again, we only need to verify that $\C$ and $H$ satisfy the conditions of \autoref{thm:koenig}. By \autoref{rem:conditions}, every extensive category satisfies (\injmono). The functor $H$ is finitary because filtered colimits commute with colimits (in particular coproducts) in all categories, and with finite limits (in particular products) in \lfp categories~\cite[Prop.~1.59]{adamek_rosicky_1994}. The functor $H$ satisfies (\hpresint) because pullbacks commute with limits (in particular products) in all categories, and with coproducts in extensive categories~\cite[Lem.~1.4]{ls05}.
\end{proofappendix}




\subsection{Nominal Transition Systems}
\label{sec:nom}
Nominal sets~\cite{pitts13} are a categorical framework to reason about variable names and
binding (like in the $\lambda$-calculus) and have also served as a basis for the 
study of automata and transition systems over infinite alphabets, whose elements represent \emph{data values}~\cite{BojanczykEA14}.

Parametric in a fixed countably infinite set of \emph{atoms} $\At$ -- intuitively the variable
names or alphabet -- a \emph{nominal set} consists of a set $X$ with
a group action $\textnormal{`\ensuremath{\cdot}'}\colon \Perm(\At)\times X\to X$
for the group $\Perm(\At)$ of finite permutations $\pi\colon \At\to \At$, satisfying the following
\emph{finite support} property: for each $x\in X$,
there is a finite set $S\subseteq \At$ such that
$
  \text{if }\pi(a) = a \text{ for all }a\in S
  \text{ then }\pi\cdot x = x.
$

Intuitively, elements of nominal sets are syntactic objects holding atoms from $\At$,
which can be renamed using finite permutations. The
finite support property means that each element in a nominal set can
hold only finitely many atoms. For example, the set $\At^*$ of finite words over~$\At$ forms a nominal set with group action $\pi\cdot (a_1\cdots a_n) = \pi(a_1)\cdots \pi(a_n)$. Similarly, the terms of the $\lambda$-calculus (modulo $\alpha$-equivalence) with variables from $\At$ form a nominal set whose group action is capture-avoiding renaming of free variables, e.g.\ $(a\, b)\cdot (\lambda a.\,a\, b) = \lambda c.\, c\, a$.



 Nominal sets form a full subcategory $\Nom$ of the category of $\Perm(\At)$-actions and equivariant maps (that is, maps preserving
the group action).
The category $\Nom$ is a \lfp boolean topos with the subobject
classifier $\Omega = 2$ being the two-element set (with trivial action).
Finitely presentable and finitely generated objects are precisely the
\emph{orbit-finite} nominal sets~\cite[Thm.~5.16]{pitts13}, that is, those nominal sets having
only finitely many elements up to renaming. For example, $\At^*$ is not orbit-finite, but the nominal set~$\At^n$ of words of fixed length $n$ is. We call a subset $M\subseteq X$ of a nominal set \emph{orbit-finite}
if $\set[\big]{\set{\pi\cdot x\mid \pi \in \Perm(\At)}\mid x\in M}$ is a finite set, and \emph{finitely supported} if it satisfies the finite support property w.r.t.\ the group action $\pi\cdot M = \set{\pi\cdot x\mid x\in M}$. The power object functor~$\Pow$ on $\Nom$ is given by $\Pow X = \set{M\seq X \mid \text{$M$ finitely supported}}$, and its finitary coreflection by $\Powf X =\set{ M\seq X \mid \text{$M$ finitely supported and orbit-finite}}$. 

Applying \autoref{thm:koenig} to the category of nominal sets gives the following general result:
\begin{theorem}[\Konig's Lemma for Coalgebras in $\Nom$]\label{thm:koenigNom}
  If $H$ is a finitary functor on $\Nom$ preserving binary intersections, then every state of a well-founded $H$-coalgebra lies in some orbit-finite
  subcoalgebra.
\end{theorem}
\begin{proofappendix}{thm:koenigNom}
  All conditions of \autoref{thm:koenig} are met because $\Nom$ is a locally finitely presentable topos, in particular extensive (cf.\ \autoref{rem:conditions})
\end{proofappendix}
Coalgebras for finitary functors on $\Nom$ admit possibly
infinite branching as long as it is finite up to renaming (cf.\ \cite[Sec.\ 6.2]{msw16}). For illustration, let us consider coalgebras for the functor $(\Pow_\omega)^\At$ on $\Nom$. (Here $(-)^\At$ denotes the exponential.) They admit a simple concrete description: A \emph{nominal labelled transition system} (\emph{NLTS}) is given by a nominal set $C$ of \emph{states} and a family of relations $\mathord{\xto{a}} \seq C\times C$  ($a\in \At$) such that $x\xto{a} y$ implies $\pi\cdot x \xto{\pi(a)} \pi\cdot y$ for all $\pi \in \Perm(\At)$. It is \emph{image-orbit-finite} if for all $x\in C$ and $a\in \At$ the set $\set{y\in C\mid x\xto{a} y}$ is orbit-finite.

An NLTS can be viewed as an ordinary LTS by forgetting the group action.

\begin{lemma}\label{lem:powfnom}
  Coalgebras for $\smash{(\Pow_\omega)^\At}$ are in bijective correspondence with image-orbit-finite NLTS. Moreover, a coalgebra is well-founded iff the
  corresponding LTS has no infinite paths.
\end{lemma}

\begin{proofappendix}{lem:powfnom}
Since $\Pow(X) = 2^X$, we have for all nominal sets $C$ and $X$  the bijective correspondence
  \[
    (C\to \Pow (X))
    \Leftrightarrow
    (C\to 2^X)
    \Leftrightarrow
    (C\times X\to 2) \Leftrightarrow \text{equivariant subsets of $C\times X$}.
  \]
Thus a coalgebra $c\colon C\to (\Pow C)^\At \cong (2^C)^\At \cong 2^{\At\times C}$ corresponds to an equivariant subset of $C\times \At\times C$. This means precisely that the corresponding family of relations $\xto{a}\,\seq C\times C$ ($a\in \At$) satifies the defining property of an NLTS. Moreover, $c$ factorizes through $(\Powf C)^\At$ iff for each $x\in C$ and $a\in \At$ the set $\set{y\in C\mid x\xto{a} y}$ is orbit-finite (by definition of $\Powf$). This means precisely that the corresponding NLTS is image-orbit-finite.

Regarding the well-foundedness statement, we can reuse the argument of \autoref{ex:wfCoalg}. We only need to note that given a coalgebra $c\colon C\to (\Powf C)^\At$, the set $S\seq C$ of all states that do not lie on any infinite path forms an equivariant subset of $C$. Since complements of equivariant subsets are equivariant, it suffices to show that $C\setminus S$ is equivariant. But this is easy to see: if $x\in C\setminus S$, there exists an infinite path $x\xto{a_1} x_1 \xto{a_2} x_2 \xto{a_3}\cdots$ starting in $x$. By equivariance of transitions, we get an infinite path $\pi\cdot x\xto{\pi(a_1)} \pi\cdot x_1 \xto{\pi(a_2)} \pi\cdot x_2 \xto{\pi(a_3)}\cdots$ for all $\pi\in\Perm(\At)$, and so $\pi\cdot x\in C\setminus S$.
\end{proofappendix}
Note that an image-orbit-finite NLTS is generally not finitely branching as an LTS, both due to the infinity of the label set $\At$ and due to the fact that each state $x$ has a possibly infinite set of $a$-successors for a given label $a$. Thus, the classical \Konig's lemma does not apply here. However, the orbit-finite branching and equivariance of transitions is enough to get following result, which is an instance of \autoref{thm:koenigNom}:

\begin{theorem}[\Konig's Lemma for Nominal LTS]\label{thm:koenigNomGraph}
  Every state of a well-founded image-orbit-finite NLTS lies in some orbit-finite subcoalgebra.
\end{theorem}
\begin{proofappendix}{thm:koenigNomGraph}
Since the functor $(\Powf)^\At$ is the composite of $\Powf$ and $(-)^\At$, is suffices to show that these two functors are finitary and preserve binary intersections (as preservation properties are maintained by composition of functors).

The functor $\Powf$ is finitary by definition and preserves binary intersections (see the proof of \autoref{thm:koenigTopos}). To see that $(-)^\At$ is finitary, we use the criterion of \autoref{lem:finitary-criterion}. We need to show that every orbit-finite equivariant subset $S\seq X^\At$ lies in $T^\At$ for some orbit-finite equivariant subset $T\seq X$. This is equivalent to saying that the corresponding equivariant map $f\colon S\times \At \to X$ factorizes through some orbit-finite equivariant subset $T\subto X$, i.e.\ its image its orbit finite. But this is clear because $S\times \At$ is orbit-finite (using that orbit-finiteness is preserved by finite products~\cite{pitts13}) and orbits are mapped to orbits by equivariant maps. Finally, the functor $(-)^\At$ preserves intersections because it is a right adjoint, thus preserves all limits.   
\end{proofappendix}

\subsection{Convex Transition Systems}
\label{sec:convex}
The category of {convex sets}~\cite{stone-49} provides a natural setting for studying probabilistic automata and transition systems, in particular systems that combine probabilistic choice with non-determinism~\cite{bss17,mv20,msv21,bsv22}. Sets of probability distributions are modelled as objects in the category, while non-determinism is captured by a convex version of the power set functor. 

A \emph{convex set} is a set \(X\) equipped with a family of binary operations \( +_{r}  \colon X \times X \rightarrow X\) (\(r \in [0, 1]\)) subject to the following equations, where \(s' = r + s -rs \ne 0\) and \(r' = \frac{r}{s'}\): 
  \begin{equation}\label{eq:convex-set-ax} x +_{r} x = x,
    \quad
    x +_{0} y = y,
    \quad
    x +_{r} y = y +_{1-r} x,
    \quad
    x +_{r} (y +_{s} z) = (x +_{r'} y ) +_{s'} z.
  \end{equation}
Convex sets are an algebraic model of structures that admit convex combinations of elements. The prototypical example of convex sets are convex subsets $X\subseteq \mathbb{R}^n$ with the operations $\vec x +_r \vec y := r\cdot \vec x
+ (1-r)\cdot \vec y$. Another example is the set $\Dist X$ of finitely supported probability distributions on a set $X$ (i.e.\ maps $\varphi\colon X\to [0,1]$ such $\varphi(x)=0$ for all but finitely many $x\in X$ and $\sum_{x\in X} \varphi(x) = 1$) with convex structure given by $\varphi+_r \psi = \big(x \mapsto r \cdot \varphi(x) + (1-r) \cdot \psi(x)\big)$.

 A map \(f \colon X \rightarrow Y\) between convex sets    is \emph{affine} if \(f(x +_{r} x') = f(x) +_{r} f(x')\) for \(x, x' \in X\) and \(r \in [0, 1]\). We let $\Conv$ denote the category of convex sets and affine maps. It is \lfp (as is every algebraic category given by a finitary equational theory~\cite{adamek_rosicky_1994}), and finitely presentable convex sets coincide with finitely generated ones~\cite[Cor.~5.5]{sw15}. Note that in algebraic categories, finitely generated objects are the algebras with a finite set of generators~\cite[Ex.~1.68]{adamek_rosicky_1994}. Thus a convex set $X$ is finitely generated iff there exists a finite subset $X_0\seq X$ such that every $x\in X$ is an (iterated) convex combination of elements of $X_0$. 

The convex analogue of the finite power set functor $\Pow_\omega$ is the \emph{convex power set functor}~$\CPowf$ on $\Conv$. It maps a convex set $X$ to the convex set $\CPowf X$ of all finitely generated (not finite!) convex subsets of~$X$. The operation $+_r$ on $\CPowf X$ is given by $S+_r T = \set{x+_r y \mid x\in S,\, y\in T}$ for $r\in (0,1)$, and determined by the second and third axiom of \eqref{eq:convex-set-ax} for $r\in \{0,1\}$ (e.g.\ $\emptyset +_0 S = S = S+_1 \emptyset$). The functor~$\CPowf$ is the composite of the \emph{non-empty} convex power set functor and the \emph{black-hole extension} functor $(-)+1$ from~\cite[Sec.~5]{bss17}.

A coalgebra $c\colon C\to \CPowf C$ is an \emph{fg-branching convex graph}, that is, a graph where (i) the set $C$ of nodes carries the structure of a convex set, (ii) for every node $x$ the set $c(x)$ of successors forms a finitely generated convex subset of $C$, and (iii) $c(x+_r y)=c(x)+_r c(y)$ for all $x,y\in C$ and $r\in [0,1]$. An fg-branching convex graph can be viewed as an ordinary graph by forgetting the convex structure on $C$. Well-foundedness is then the expected concept:

\begin{lemma}\label{lem:convex-well-founded}
A coalgebra $c\colon C\to \CPowf C$ is well-founded iff it has no infinite paths.
\end{lemma}

\begin{proofappendix}{lem:convex-well-founded}
The argument is analogous to the case of graphs (\autoref{ex:wfCoalg}). The only thing we need to additionally observe is that given a coalgebra $c\colon C\to \CPowf C$, the set $S\seq C$ of all nodes that lie on no infinite path forms a convex subset of $C$. Thus let $x,y\in S$ and $z=x+_r y$ for some $r\in [0,1]$. We need to prove $z\in S$. If $r\in \{0,1\}$ then $z\in \{x,y\}$ and so $z\in S$. Now let $r\in (0,1)$, and suppose towards a contradiction that $z\not\in S$, that is, there is an infinite path $z\to z_1 \to z_2\to\ldots$ starting with $z$. Then $z_1\in c(z)=c(x+_r y) = c(x)+_r c(y)$. This means that there exist $x\to x_1$ and $y\to y_1$ such $z_1=x_1+_r y_1$. Repeating this argument with $x_1,y_1,z_1$ instead of $x,y,z$ yields $x_1\to x_2$ and $y_1\to y_2$ such that $z_2=x_2+_r y_2$ etc. In this way, we obtain an infinite path $x\to x_1\to x_2\to \cdots$, in  contradiction to $x\not\in S$.    
\end{proofappendix}



When regarded as an ordinary graph, an fg-branching convex graph is generally not finitely branching. However, the fg-branching property is enough to get a \Konig's lemma. Indeed, applying \autoref{thm:koenig} to the category $\Conv$ and the functor $\CPowf$ yields:

\begin{theorem}[\Konig's Lemma for Convex Graphs]\label{thm:koenig-convex}
Every node of a well-founded fg-branching convex graph lies in some fg-carried $\CPowf$-subcoalgebra.
\end{theorem}

\begin{proofappendix}{thm:koenig-convex}
As usual, we only need to show that the conditions of \autoref{thm:koenig} are satisfied. The condition (\injmono) follows from the description of coproducts in $\Conv$ given in~\cite[Prop.~5]{jww15}. Clearly, the functor $\CPowf$ preserves binary intersections. (Note that since $\Conv$ is an algebraic category, limits in $\Conv$ are formed like in $\Set$.) Finally, to check that $\CPowf$ is finitary, we use the criterion of \autoref{lem:finitary-criterion}. Thus, we need to show that every finitely generated convex set $\mathcal{S}\seq \CPowf C$ is a subset of $\CPowf M$ for some finitely generated convex set $M\seq C$. But this is easy: Let $\S_0$ be a finite set of generators of $\S$. Then each $S\in \S_0$ is a finitely generated convex subset of $C$ and thus has a finite set of generators $S_0\seq S$. Let $M\seq X$ be the convex subset generated by the finite set $\bigcup_{S\in \S_0} S_0$. Then $S\seq M$ for every $S\in \S_0$, and so $S\seq M$ for every $S\in \S$ because $\S$ is generated by $\S_0$.
\end{proofappendix}

\section{Initial Algebras From Finite Well-Founded and Recursive Coalgebras}\label{sec:initial-algebras}
In this section we establish our second fundamental result about well-founded coalgebras: the initial algebra for $H$ can be constructed as the colimit of all fp-carried well-founded coalgebras. Note that, in contrast, \autoref{thm:koenig} involves \emph{fg-carried} well-founded coalgebras. The proof is another intricate application of coproduct extension (\autoref{cons:cp}).

\begin{construction}\label{not:T}
Let $(T,t)$ be the colimit of all coalgebras in $\fwf$, i.e.\ the colimit of the embedding
$D\colon \Coalg_{\fp,\wf}(H)\subto \Coalg(H)$.
We denote the colimit injections by
\[ c^\#\colon (C,c)\to (T,t) \qquad \text{where} \qquad (C,c)\in \fwf. \]
This colimit is formed at the level of the underlying category $\C$ (\autoref{lem:colimits-coalg}). Moreover, if $\C$  has a simple initial object, e.g.~if (\injmono) holds (\autoref{rem:conditions}), the colimit is filtered because the category $\fwf$ is finitely cocomplete. Indeed, both $\C_\fp\hookrightarrow \C$ and $\Coalg_\wf(H)\hookrightarrow \Coalg(H)$ are closed under finite colimits (\autoref{lem:fp-fg-coalg}, \autoref{prop:wf}).
\end{construction}
We then have the following characterization of the initial algebra for $H$:

\begin{theorem}[First Initial Algebra Theorem,\,\smash{\injmono},\,\hpresint]\label{thm:wf-initial-algebra}
The initial algebra for $H$ is the colimit of all fp-carried well-founded coalgebras, that is, the coalgebra structure $t\colon T\to HT$ is an isomorphism and the algebra $(T,t^{-1})$ is initial for $H$.
\end{theorem}

The crucial part of this theorem is the `Lambek lemma' for the coalgebra $(T,t)$:
\begin{lemma}[\smash{\injmono}, \hpresint]\label{lem:t-iso-wf}
The coalgebra structure $t\colon T\to HT$ is an isomorphism.
\end{lemma}

\begin{proof}[Proof sketch]
Let $U\colon \Coalg(H)\to \C$ be the forgetful functor. Then $(c^\#)$ is a colimit cocone over $UD$ and $(Hc^\#\circ c)$ is a cocone over $UD$, and $t$ is the unique mediating morphism such that the diagram below commutes for every $(C,c)\in \fwf$. It suffices to prove that
  $(Hc^\#\circ c)$ forms a colimit cocone over $UD$; then $t$ is an isomorphism by uniqueness of colimits. To this end, we verify the criterion of \autoref{lem:filtered-colimits}: for $X\in \C_\fp$, every morphism $q\colon X \to HT$ factorizes essentially uniquely through the cocone $(Hc^\# \circ c)$.
\[  
\begin{tikzcd}[row sep=3mm,outer sep=0pt]
T \ar[dashed]{r}{t} & HT \\
C \ar{u}[outer sep=5pt]{c^\#} \ar{r}{c} & HC \ar{u}[swap,outer sep=5pt]{Hc^\#}
\end{tikzcd}
\]
To construct the factorization, we first observe that since $(c^\#)$ is a filtered colimit and $H$ is finitary, $(Hc^\#)$ is a filtered colimit. Therefore, since $X\in \C_\fp$, the morphism $q$ factorizes as $q=Hc^\# \circ p$ for some $(C,c)\in \Coalg_{\fp,\wf}(H)$ and $p\colon X\to HC$. Consider the coproduct extension $(C+X,c_p)$ of $(C,c)$ by $p$ (\autoref{cons:cp}). Note that $(C+X,c_p)\in \Coalg_{\fp,\wf}(H)$ because coproduct extension preserves well-foundedness (\autoref{prop:wf-coalgebra-coproduct}) and $\C_\fp\subto \C$ is closed under finite coproducts (\autoref{lem:fp-fp-fin-colimits}). It is not difficult to see that $q\colon X\to HT$ satisfies $q=(\,X \flatxto{\inr} C+X \flatxto{c_p}  H(C+X)\flatxto{Hc_p^\#}  HT\,)$; hence $q$ factorizes through $(Hc^\# \circ c)$.
\end{proof}

\begin{proofappendix}{lem:t-iso-wf}
Let $U\colon \Coalg(H)\to \C$ denote the forgetful functor. Then $(c^\#)$ is a colimit cocone over $UD$ and $(Hc^\#\circ c)$ is a cocone over $UD$, and $t$ is the unique mediating morphism such that
\[  
\begin{tikzcd}
T \ar[dashed]{r}{t} & HT \\
C \ar{u}{c^\#} \ar{r}{c} & HC \ar{u}[swap]{Hc^\#}
\end{tikzcd}
\]
commutes for every $(C,c)\in \fwf$. (This observation underlies the proof that $U$ creates colimits, see \autoref{lem:colimits-coalg}.) It suffices to show that $(Hc^\#\circ c)$ forms a colimit cocone over $UD$; then $t$ is an isomorphism by the uniqueness of colimits. To this end, we verify the criterion of \autoref{lem:filtered-colimits}, i.e.\ we show that for $X\in \C_\fp$, every morphism from $X$ to $HT$ factorizes essentially uniquely through the cocone $(Hc^\# \circ c)$.
\begin{enumerate}
\item Let $q\colon X\to HT$ be a morphism. Since $H$ is finitary, $(Hc^\#)$ is a colimit cocone over $HUD$. Hence there exists $(C,c)\in \fwf$  and $p\colon X\to HC$ such that
\[ q=Hc^\# \circ p. \]
Consider the coproduct extension $(C+X,c_p)$ of $(C,c)$ by $p$ (\autoref{cons:cp}). By \autoref{prop:wf-coalgebra-coproduct} and \autoref{lem:fp-fp-fin-colimits}, we have $(C+X,c_p)\in \fwf$. Moreover, the following diagram commutes:
\[
\begin{tikzcd}
X \ar{dd}[swap]{\inr} \ar{rr}{q} \ar{dr}{p} & & HT \\
& HC \ar{ur}{Hc^\#} \ar{dr}{H\inl} & \\
C+X \ar{ur}{[c,p]} \ar{rr}{c_p} & & H(C+X) \ar{uu}[swap]{Hc_p^\#}
\end{tikzcd}
\]
Indeed, the right-hand triangle commutes because $\inl\colon (C,c)\to (C+X,c_p)$ is a coalgebra morphism (\autoref{lem:coproduct-inl-coalg-mor}) and $(-)^\#$ is a cocone. The remaining parts clearly commute. Hence the outside commutes, which proves that $q$ factorizes through $Hc_p^\# \circ c_p$ via $\inr$.
\item To prove essential uniqueness of factorizations, let $(C,c)\in \fwf$ and let $f,f'\colon X\to C$ be morphisms such that 
\[ q:= Hc^\# \circ c \circ f = Hc^\# \circ c \circ f'.\]
Our task is to produce a coalgebra morphism in $\fwf$ that merges $f$ and $f'$. Put
\[ p:=c \circ f\qqand p':=c\circ f'.  \]
 Since $(Hc^\#)$ is a colimit cocone over $HUD$ and $X\in \C_\fp$, there exists $g\colon (C,c)\to (\ol{C},\ol{c})$ in $\fwf$ such that  
\[ Hg\circ p = Hg\circ p'.\]
Therefore, after replacing $f,f'$ with $g\circ f, g\circ f'$, we may assume w.l.o.g.\ that $p=p'$.
\[
\begin{tikzcd}[row sep=3em, column sep=5em]
X \ar{r}{q} \ar[shift right=1]{dr}[swap]{p} \ar[shift left=1]{dr}{p'} \ar[shift right=1, bend right=10]{ddr}[swap]{f} \ar[shift left=1, bend right=10]{ddr}{f'} & HT & \\
& HC \ar{r}{Hg} \ar{u}[swap]{Hc^\#} & H\ol{C} \ar{ul}[swap]{H\ol{c}^\#} \\
& C \ar{u}[swap]{c} \ar{r}{g} & \ol{C} \ar{u}[swap]{\ol{c}}
\end{tikzcd}
\]
Then $[\id,f],[\id,f']\colon (C+X,c_p)\to (C,c)$ are coalgebra morphisms; this is shown by the commutative diagram below for $[\id,f]$, and analogously for $[\id,f']$, using that $p=p'$.
\[
\begin{tikzcd}
C+X \ar{r}{[\id,f]}  \ar{d}[swap]{[c,p]} \ar[shiftarr={xshift=-14mm}]{dd}[swap]{c_p} & C \ar{dd}{c} \\
HC \ar{d}[swap]{H\inl} & \\
H(C+X) \ar{r}{H[\id,f]} & HC
\end{tikzcd}
\]
Since $\fwf$ is filtered, there exists a morphism  $h\colon (C,c)\to (C',c')$ in $\fwf$ merging $[\id,f]$ and $[\id,f']$:
\[  
\begin{tikzcd}
(C+X,c_p) \ar[shift left=1]{r}{[\id,f]} \ar[shift right=1]{r}[swap]{[\id,f']} & (C,c) \ar{r}{h} & (C',c').
\end{tikzcd}
\]
In particular, we have $h\circ f= h\circ f'$, as required.\qedhere
\end{enumerate}
\end{proofappendix}

From the `Lambek lemma', the First Initial Algebra Theorem is immediate:


\begin{proof}[Proof of \autoref{thm:wf-initial-algebra}]
The coalgebra $(T,t)$ is well-founded because it is a colimit of well-founded coalgebras (\iref{prop:wf}{closed-colim-quotients}), and so it is recursive (\iref{prop:wf}{wf-to-rec}). Clearly, every recursive coalgebra whose structure is an isomorphism is initial as an algebra.
\end{proof}

Wißmann and Milius~\cite{wm24} have recently shown an analogue of \autoref{thm:wf-initial-algebra} for \emph{recursive} instead of well-founded coalgebras. Consider the following modification of \autoref{not:T}:

\begin{construction}\label{not:barT}
Let $(\ol{T},\ol{t})$ be the colimit of all coalgebras in $\frc$, i.e.\ the colimit of the embedding $\ol{D}\colon \frc\subto \Coalg(H)$. 
\end{construction}
This colimit is filtered because $\Coalg_{\fp,\rec}(H)$ is finitely cocomplete, using that both $\C_\fp\subto \C$ and $\Coalg_\rec(H)\subto \Coalg(H)$ are closed under finite colimits (\autoref{lem:fp-fp-fin-colimits}, \autoref{lem:rec-coalg-colimits}). It turns out that \autoref{not:barT} also yields the initial algebra:

\begin{theorem}[Second Initial Algebra Theorem~\cite{wm24}]\label{thm:rec-initial-algebras}
The initial algebra for $H$ is the colimit of all fp-carried recursive coalgebras, that is, the coalgebra structure $\ol{t}\colon \ol{T}\to H\ol{T}$ is an isomorphism and the algebra $(\ol{T},\ol{t}^{-1})$ is initial for $H$.
\end{theorem}
We stress that \autoref{thm:rec-initial-algebras} does not need any conditions on $\C$ and $H$ beyond those of \autoref{asm}. However, in the case where the additional conditions (\injmono) and (\hpresint) are satisfied, the first initial algebra theorem can be regarded as an improvement of the second one, as it constructs the initial algebra from a smaller diagram. Indeed, under these conditions well-founded coalgebras form a subclass of recursive ones (\autoref{prop:wf}), and this subclass is generally proper, as witnessed by the coalgebra in \autoref{ex:konig-lemma-recursive}.

The idea underlying our proof of \autoref{thm:wf-initial-algebra} carries over easily to the recursive case and leads to a new proof of \autoref{thm:rec-initial-algebras} that is substantially shorter and simpler than the original one~\cite{wm24}. Essentially, the only thing we need to observe is that coproduct extension not only preserves well-foundedness (\autoref{prop:wf-coalgebra-coproduct}), but also recursivity:

\begin{proposition}\label{prop:rec-coalgebra-coproduct}
For every recursive coalgebra $(C,c)$ and every morphism $p\colon X\to HC$, the coproduct extension $(C+X,c_p)$ is a recursive coalgebra.
\end{proposition}

\begin{proofappendix}{prop:rec-coalgebra-coproduct}
Suppose that $(C,c)$ is a recursive coalgebra, and let $(A,a)$ be an $H$-algebra. We need to show that there exists a unique coalgebra-to-algebra morphism from $(C+X,c_p)$ to $(A,a)$.

\medskip\noindent\emph{Existence.} Since $(C,c)$ is recursive, there exists a unique coalgebra-to-algebra morphism $h\colon (C,c)\to (A,a)$. We claim that the composite
\[ C+X \xto{\id+p} C+HC \xto{h+Hh} A+HA \xto{[\id,a]} A \]
is a coalgebra-to-algebra morphism from $(C+X, c_p)$ to $(A,a)$, which means that the outside of the diagram below commutes:
\[
\begin{tikzcd}[column sep=35]
C+X \ar[shiftarr={xshift=-14mm}]{dd}[swap]{c_p} \ar{r}{\id+p} \ar{d}[swap]{[c,p]} & C+HC \ar{r}{h+Hh} & A+HA \ar{r}{[\id,a]} & A \\
HC \ar{rrr}{Hh} \ar{d}[swap]{H\inl} & & & HA \ar{u}[swap]{a} \\
H(C+X) \ar{r}{H(\id+p)} & H(C+HC) \ar{r}{H(h+Hh)} & H(A+HA) \ar{r}{H([\id,a])} & HA \ar[equals]{u}
\end{tikzcd}
\]
Indeed, the lower part commutes trivially, and the left-hand part by definition of $c_p$. For the upper part consider its precomposition with the coproduct injections of $C+X$. The precomposition with $\inl$ yields the diagram expressing that $h$ is a coalgebra-to-algebra morphism, and the precomposition with $\inr$ commutes trivially.

\medskip\noindent\emph{Uniqueness.} Let $[h,k]\colon (C+X,c_p)\to (A,a)$ be a coalgebra-to-algebra morphism. Then the outside of the diagram
\[
\begin{tikzcd}
C+X \ar[shiftarr={xshift=-18mm}]{dd}[swap]{c_p} \ar{r}{[h,k]} \ar{d}[swap]{[c,p]} & A \\
HC \ar{r}{Hh} \ar{d}[swap]{H\inl} & HA \ar{u}[swap]{a} \\
H(C+X) \ar{r}{H[h,k]} & HA \ar[equals]{u}
\end{tikzcd}
\]
commutes, as do its lower and left-hand parts. Hence the upper part commutes. Precomposing it with $\inl$ shows that $h\c (C,c)\to (A,a)$ is a coalgebra-to-algebra morphism; thus $h$ is uniquely determined because $(C,c)$ is recursive. Precomposition with $\inr$ shows that $k=a\circ Hh\circ p$, so $k$ is determined by $h$. \qedhere 
\end{proofappendix}
Using this result, we can now proceed exactly like in the proof of \autoref{thm:wf-initial-algebra}. First, we establish the `Lambek lemma' for the coalgebra $(\ol{T},\ol{t})$:

\begin{lemma}\label{lem:t-iso}
The coalgebra structure $\ol{t}\colon \ol{T}\to H\ol{T}$ is an isomorphism.
\end{lemma}
The proof is identical to that of \autoref{lem:t-iso-wf}, except than one considers $\frc$ and $\ol{D}$ instead of $\fwf$ and $D$, and uses \autoref{prop:rec-coalgebra-coproduct} to show $(C+X,c_p)\in \frc$. 

From the `Lambek lemma', the Second Initial Algebra Theorem is now again immediate:

\begin{proof}[Proof of \autoref{thm:rec-initial-algebras}]
The coalgebra $(\ol{T},\ol{t})$ is recursive because it is a colimit of recursive coalgebras (\autoref{lem:rec-coalg-colimits}). Therefore $(\ol{T},\ol{t}^{-1})$ is an initial algebra.
\end{proof}

\begin{rem}
It is instructive to compare \autoref{thm:wf-initial-algebra} and \ref{thm:rec-initial-algebras} with the well-known iterative construction of initial algebras due to  Ad\'amek~\cite{adamek74}:
For every finitary endofunctor~$H$ on a cocomplete category $\C$, the initial algebra is the colimit of the $\omega$-chain
\[
(0,\initial)\xto{\smash{\initial}}
(H0,H\initial)\xto{\smash{H\initial}}
(HH0,HH\initial) \xto{\smash{HH\initial}}
\cdots \xto{\smash{H^{n-1}\initial}} 
(H^n0,H^n\initial)\xto{\smash{H^n\initial}} \cdots 
\]
in $\Coalg(H)$, where $0$ is the initial object of $\C$ and $\initial\colon 0\to H0$ is the unique morphism. Similar to \autoref{thm:wf-initial-algebra} and \ref{thm:rec-initial-algebras}, this construction approximates the initial algebra from below: intuitively, one thinks of the initial algebra as the set of all finite trees with branching type $H$, and of $H^n 0$ as the set of trees of height less than $n$. The approximating coalgebras $(H^n0,H^n\initial)$ are recursive~\cite[Prop.~6]{cuv06}, but in contrast to \autoref{thm:wf-initial-algebra} and \ref{thm:rec-initial-algebras}, they are generally neither well-founded nor fp-carried.
\end{rem}


\section{Conclusion and Future Work}
We have established two novel results in the theory of well-founded coalgebras over locally finitely presentable categories. Our first contribution is a coalgebraic version of \Konig's lemma, lifting a previous result for coalgebras over sets~\cite{amm25} to full categorical generality. We have illustrated the scope of our coalgebraic \Konig's lemma by exploring systems with state spaces beyond sets, such as graphs in a topos, and nominal and convex transition systems. As our second contribution, we devised a new construction of the initial algebra for a finitary functor as the colimit of all well-founded coalgebras with finitely presentable state space. 

One key application of well-foundedness is termination analysis of programs~\cite{cpr11}, which amounts to associating cleverly chosen well-founded relations to the state graph of a program. We aim to explore a coalgebraic perspective on termination analysis techniques, in particular in relation with the recent advances in modelling stateful languages as coalgebras~\cite{gmstu25}. 

A further interesting route is studying connections between well-founded coalgebras and \emph{well-structured transition systems}~\cite{fs01}, which involve a well-quasi-order on states and provide a uniform framework for a class of decidability results in verification, e.g.~for Petri nets.




\providecommand{\noopsort}[1]{}
\begin{thebibliography}{10}

\bibitem{ap04}
J.~Ad\'amek and H.-E. Porst.
\newblock On tree coalgebras and coalgebra presentations.
\newblock {\em Theoretical Computer Science}, 311(1):257--283, 2004.
\newblock \href {https://doi.org/10.1016/S0304-3975(03)00378-5}
  {\path{doi:10.1016/S0304-3975(03)00378-5}}.

\bibitem{amv03}
Jir\'{\i} Ad\'{a}mek, Stefan Milius, and Jir\'{\i} Velebil.
\newblock Free iterative theories: a coalgebraic view.
\newblock {\em Math.~Struct.~Comput.~Sci.}, 13(2):259–320, 2003.
\newblock \href {https://doi.org/10.1017/S0960129502003924}
  {\path{doi:10.1017/S0960129502003924}}.

\bibitem{adamek74}
Ji\v{r}\'i Ad\'amek.
\newblock Free algebras and automata realizations in the language of
  categories.
\newblock {\em Comment. Math. Univ. Carol.}, 15(4):589--602, 1974.

\bibitem{amm25}
Ji\v{r}\'i Ad\'amek, Stefan Milius, and Lawrence~S. Moss.
\newblock {\em Initial Algebras and Terminal Coalgebras: The Theory of Fixed
  Points of Functors}.
\newblock Cambridge Tracts in Theoretical Computer Science. Cambridge
  University Press, 2025.
\newblock \href {https://doi.org/10.1017/9781108884112}
  {\path{doi:10.1017/9781108884112}}.

\bibitem{amsw19}
Ji\v{r}\'i{} Ad\'amek, Stefan Milius, Lurdes Sousa, and Thorsten Wi{\ss}mann.
\newblock On finitary functors.
\newblock {\em Theory Appl. Categ.}, 34(35):1134--1164, 10 2019.

\bibitem{amv06}
Ji\v{r}\'i Ad\'amek, Stefan Milius, and Ji\v{r}\'i Velebil.
\newblock Iterative algebras at work.
\newblock {\em Math.~Struct.~Comput.~Sci.}, 16(6):1085–1131, 2006.
\newblock \href {https://doi.org/10.1017/S0960129506005706}
  {\path{doi:10.1017/S0960129506005706}}.

\bibitem{adamek_rosicky_1994}
Ji\v{r}\'i Ad\'amek and Ji\v{r}\'i Rosick\'y.
\newblock {\em Locally Presentable and Accessible Categories}.
\newblock London Mathematical Society Lecture Note Series. Cambridge University
  Press, 1994.
\newblock \href {https://doi.org/10.1017/CBO9780511600579}
  {\path{doi:10.1017/CBO9780511600579}}.

\bibitem{BojanczykEA14}
Miko{\l}aj Boja{\'n}czyk, Bartek Klin, and S{\l}awomir Lasota.
\newblock Automata theory in nominal sets.
\newblock {\em Log. Methods Comput. Sci.}, 10(3), 2014.
\newblock \href {https://doi.org/10.2168/LMCS-10(3:4)2014}
  {\path{doi:10.2168/LMCS-10(3:4)2014}}.

\bibitem{bss17}
Filippo Bonchi, Alexandra Silva, and Ana Sokolova.
\newblock The power of convex algebras.
\newblock In {\em 28th International Conference on Concurrency Theory (CONCUR
  2017)}, volume~85 of {\em LIPIcs}, pages 23:1--23:18. Schloss Dagstuhl --
  Leibniz-Zentrum f{\"u}r Informatik, 2017.
\newblock \href {https://doi.org/10.4230/LIPIcs.CONCUR.2017.23}
  {\path{doi:10.4230/LIPIcs.CONCUR.2017.23}}.

\bibitem{bsv22}
Filippo Bonchi, Ana Sokolova, and Valeria Vignudelli.
\newblock The theory of traces for systems with nondeterminism, probability,
  and termination.
\newblock {\em Log. Methods Comput. Sci.}, 18(2), 2022.
\newblock \href {https://doi.org/10.46298/lmcs-18(2:21)2022}
  {\path{doi:10.46298/lmcs-18(2:21)2022}}.

\bibitem{brouwer27}
Luitzen E.~J. Brouwer.
\newblock {\"U}ber {D}efinitionsbereiche von {F}unktionen.
\newblock {\em Math.~Ann.}, 97(1):60--75, 1927.
\newblock \href {https://doi.org/10.1007/BF01447860}
  {\path{doi:10.1007/BF01447860}}.

\bibitem{brv04}
Doron Bustan, Sasha Rubin, and Moshe~Y. Vardi.
\newblock Verifying $\omega$-regular properties of {M}arkov chains.
\newblock In {\em 16th International on Computer Aided Verification (CAV
  2004)}, volume 3114 of {\em Lect.~Notes~Comput.~Sci.}, pages 189--201.
  Springer, 2004.
\newblock \href {https://doi.org/10.1007/978-3-540-27813-9_15}
  {\path{doi:10.1007/978-3-540-27813-9_15}}.

\bibitem{cuv06}
Venanzio Capretta, Tarmo Uustalu, and Varmo Vene.
\newblock Recursive coalgebras from comonads.
\newblock {\em Inf. Comput.}, 204(4):437--468, 2006.
\newblock \href {https://doi.org/10.1016/J.IC.2005.08.005}
  {\path{doi:10.1016/J.IC.2005.08.005}}.

\bibitem{cw93}
Aurelio Carboni, Stephen Lack, and R.F.C. Walters.
\newblock Introduction to extensive and distributive categories.
\newblock {\em J.~Pure~Appl.~Algebra}, 84(2):145--158, 1993.
\newblock \href {https://doi.org/10.1016/0022-4049(93)90035-R}
  {\path{doi:10.1016/0022-4049(93)90035-R}}.

\bibitem{cpr11}
Byron Cook, Andreas Podelski, and Andrey Rybalchenko.
\newblock Proving program termination.
\newblock {\em Commun.~ACM}, 54(5):88–98, 2011.
\newblock \href {https://doi.org/10.1145/1941487.1941509}
  {\path{doi:10.1145/1941487.1941509}}.

\bibitem{di21}
Hannes Diener and Hajime Ishihara.
\newblock Bishop-style constructive reverse mathematics.
\newblock In Vasco Brattka and Peter Hertling, editors, {\em Handbook of
  Computability and Complexity in Analysis}, pages 347--365. Springer, 2021.
\newblock \href {https://doi.org/10.1007/978-3-030-59234-9_10}
  {\path{doi:10.1007/978-3-030-59234-9_10}}.

\bibitem{fs01}
A.~Finkel and Ph. Schnoebelen.
\newblock Well-structured transition systems everywhere!
\newblock {\em Theor.~Comput.~Sci.}, 256(1):63--92, 2001.
\newblock \href {https://doi.org/10.1016/S0304-3975(00)00102-X}
  {\path{doi:10.1016/S0304-3975(00)00102-X}}.

\bibitem{gmms14}
Tomasz Gogacz, Henryk Michalewski, Matteo Mio, and Micha{\l} Skrzypczak.
\newblock Measure properties of game tree languages.
\newblock In {\em 39th International Symposium on Mathematical Foundations of
  Computer Science (MFCS 2014)}, volume 8635 of {\em Lect.~Notes~Comput.~Sci.},
  pages 303--314. Springer, 2014.
\newblock \href {https://doi.org/10.1007/978-3-662-44522-8_26}
  {\path{doi:10.1007/978-3-662-44522-8_26}}.

\bibitem{gmstu25}
Sergey Goncharov, Stefan Milius, Lutz Schröder, Stelios Tsampas, and Henning
  Urbat.
\newblock Bialgebraic reasoning on stateful languages, 2025.
\newblock To appear in 30th ACM SIGPLAN International Conference on Functional
  Programming (IFCP 2025).
\newblock \href {https://arxiv.org/abs/2503.10955} {\path{arXiv:2503.10955}}.

\bibitem{hirschfeldt14}
Denis~R. Hirschfeldt.
\newblock {\em Slicing the Truth: On the Computable and Reverse Mathematics of
  Combinatorial Principles}.
\newblock World Scientific, 2014.

\bibitem{hr98}
Paul Howard and Jean~E. Rubin.
\newblock {\em Consequences of the Axiom of Choice}.
\newblock American Mathematical Society, 1998.
\newblock \href {https://doi.org/10.1017/S1079898600003449}
  {\path{doi:10.1017/S1079898600003449}}.

\bibitem{isb16}
Tomasz Idziaszek, Micha{\l} Skrzypczak, and Miko{\l}aj Boja{\'n}czyk.
\newblock Regular languages of thin trees.
\newblock {\em Theory Comput. Syst.}, 58(4):614--663, 2016.
\newblock \href {https://doi.org/10.1007/S00224-014-9595-Z}
  {\path{doi:10.1007/S00224-014-9595-Z}}.

\bibitem{jacobs16}
Bart Jacobs.
\newblock {\em Introduction to Coalgebra: Towards Mathematics of States and
  Observation}.
\newblock Cambridge Tracts in Theoretical Computer Science. Cambridge
  University Press, 2016.
\newblock \href {https://doi.org/10.1017/CBO9781316823187}
  {\path{doi:10.1017/CBO9781316823187}}.

\bibitem{jww15}
Bart Jacobs, Bas Westerbaan, and Bram Westerbaan.
\newblock States of convex sets.
\newblock In {\em 18th International Conference on Foundations of Software
  Science and Computation Structures (FOSSACS 2015)}, volume 9034 of {\em
  Lect.~Notes~Comput.~Sci.}, pages 87--101. Springer, 2015.
\newblock \href {https://doi.org/10.1007/978-3-662-46678-0_6}
  {\path{doi:10.1007/978-3-662-46678-0_6}}.

\bibitem{jech02}
Thomas Jech.
\newblock {\em Set Theory: The Third Millennium Edition}.
\newblock Springer, 2002.
\newblock \href {https://doi.org/10.1007/3-540-44761-X}
  {\path{doi:10.1007/3-540-44761-X}}.

\bibitem{Johnstone}
Peter~T Johnstone.
\newblock {\em {Sketches of an elephant}}.
\newblock Oxford logic guides. Oxford Univ. Press, New York, NY, 2002.
\newblock \href {https://doi.org/10.1017/S1079898600003462}
  {\path{doi:10.1017/S1079898600003462}}.

\bibitem{Johnstone_short}
Peter~T Johnstone.
\newblock {\em {Topos Theory}}.
\newblock Dover Publications, 2014.

\bibitem{KaminskiFrancez94}
Michael Kaminski and Nissim Francez.
\newblock Finite-memory automata.
\newblock {\em Theor.~Comput.~Sci.}, 134(2):329--363, 1994.
\newblock \href {https://doi.org/10.1016/0304-3975(94)90242-9}
  {\path{doi:10.1016/0304-3975(94)90242-9}}.

\bibitem{konig1927}
D{\'e}nes K\H{o}nig.
\newblock {\"U}ber eine {S}chlussweise aus dem {E}ndlichen ins {U}nendliche.
\newblock {\em Acta Sci. Math. (Szeged)}, 3(2-3):121--130, 1927.

\bibitem{ls05}
Stephen Lack and Pawel Sobocinski.
\newblock Adhesive and quasiadhesive categories.
\newblock {\em {RAIRO} Theor. Informatics Appl.}, 39(3):511--545, 2005.
\newblock \href {https://doi.org/10.1051/ITA:2005028}
  {\path{doi:10.1051/ITA:2005028}}.

\bibitem{lambek68}
Joachim Lambek.
\newblock A fixpoint theorem for complete categories.
\newblock {\em Mathematische Zeitschrift}, 103:151--161, 1968.
\newblock \href {https://doi.org/10.1007/BF01110627}
  {\path{doi:10.1007/BF01110627}}.

\bibitem{mac2013categories}
Saunders Mac~Lane.
\newblock {\em Categories for the {{Working Mathematician}}}, volume~5 of {\em
  Graduate {{Texts}} in {{Mathematics}}}.
\newblock Springer, 2 edition, 1978.
\newblock \href {https://doi.org/10.1007/978-1-4757-4721-8}
  {\path{doi:10.1007/978-1-4757-4721-8}}.

\bibitem{mpw16}
Stefan Milius, Dirk Pattinson, and Thorsten Wi{\ss}mann.
\newblock A new foundation for finitary corecursion -- the locally finite
  fixpoint and its properties.
\newblock In Bart Jacobs and Christof L{\"{o}}ding, editors, {\em Proc.~19th
  International Conference on Foundations of Software Science and Computation
  Structures (FoSSaCS 2016)}, volume 9634 of {\em LNCS}, pages 107--125.
  Springer, 04 2016.
\newblock \href {https://doi.org/10.1007/978-3-662-49630-5_7}
  {\path{doi:10.1007/978-3-662-49630-5_7}}.

\bibitem{lffiac}
Stefan Milius, Dirk Pattinson, and Thorsten Wi{\ss}mann.
\newblock A new foundation for finitary corecursion and iterative algebras.
\newblock {\em Information and Computation}, 271:1--28, 04 2020.
\newblock Article 104456.
\newblock \href {https://doi.org/10.1016/j.ic.2019.104456}
  {\path{doi:10.1016/j.ic.2019.104456}}.

\bibitem{msw16}
Stefan Milius, Lutz Schr{\"{o}}der, and Thorsten Wi{\ss}mann.
\newblock Regular behaviours with names.
\newblock {\em Applied Categorical Structures}, 24(5):663--701, 08 2016.
\newblock \href {https://doi.org/10.1007/s10485-016-9457-8}
  {\path{doi:10.1007/s10485-016-9457-8}}.

\bibitem{msv21}
Matteo Mio, Ralph Sarkis, and Valeria Vignudelli.
\newblock Combining nondeterminism, probability, and termination: Equational
  and metric reasoning.
\newblock In {\em 36th Annual ACM/IEEE Symposium on Logic in Computer Science
  (LICS 2021)}, pages 1--14, 2021.
\newblock \href {https://doi.org/10.1109/LICS52264.2021.9470717}
  {\path{doi:10.1109/LICS52264.2021.9470717}}.

\bibitem{mv20}
Matteo Mio and Valeria Vignudelli.
\newblock {Monads and Quantitative Equational Theories for Nondeterminism and
  Probability}.
\newblock In {\em 31st International Conference on Concurrency Theory (CONCUR
  2020)}, volume 171 of {\em LIPIcs}, pages 28:1--28:18. Schloss Dagstuhl --
  Leibniz-Zentrum f{\"u}r Informatik, 2020.
\newblock \href {https://doi.org/10.4230/LIPIcs.CONCUR.2020.28}
  {\path{doi:10.4230/LIPIcs.CONCUR.2020.28}}.

\bibitem{osius74}
Gerhard Osius.
\newblock Categorical set theory: A characterization of the category of sets.
\newblock {\em J.~Pure~Appl.~Algebra}, 4(1):79--119, 1974.
\newblock \href {https://doi.org/10.1016/0022-4049(74)90032-2}
  {\path{doi:10.1016/0022-4049(74)90032-2}}.

\bibitem{pitts13}
Andrew~M. Pitts.
\newblock {\em Nominal Sets: Names and Symmetry in Computer Science}.
\newblock Cambridge Tracts in Theoretical Computer Science. Cambridge
  University Press, 2013.
\newblock \href {https://doi.org/10.1017/CBO9781139084673}
  {\path{doi:10.1017/CBO9781139084673}}.

\bibitem{rutten00}
Jan J. M.~M. Rutten.
\newblock Universal coalgebra: a theory of systems.
\newblock {\em Theor. Comput. Sci.}, 249(1):3--80, 2000.
\newblock \href {https://doi.org/10.1016/S0304-3975(00)00056-6}
  {\path{doi:10.1016/S0304-3975(00)00056-6}}.

\bibitem{sw15}
Ana Sokolova and Harald Woracek.
\newblock Congruences of convex algebras.
\newblock {\em J.~Pure~Appl.~Algebra}, 219(8):3110--3148, 2015.
\newblock \href {https://doi.org/10.1016/j.jpaa.2014.10.005}
  {\path{doi:10.1016/j.jpaa.2014.10.005}}.

\bibitem{stone-49}
Marshall~H. Stone.
\newblock Postulates for the barycentric calculus.
\newblock {\em Ann. Mat. Pura Appl.}, 29(1):25--30, 1949.
\newblock \href {https://doi.org/10.1007/BF02413910}
  {\path{doi:10.1007/BF02413910}}.

\bibitem{szekely23}
L\'aszl\'o Sz\'ekely.
\newblock Ramsey theory ({L}ecture notes).
\newblock \url{https://people.math.sc.edu/laszlo/Ramsey2.pdf}, 2023.

\bibitem{taylor99}
Paul Taylor.
\newblock {\em Practical Foundations of Mathematics}.
\newblock Cambridge Studies in Advanced Mathematics. Cambridge University
  Press, 1999.
\newblock \href {https://doi.org/10.2307/3621547} {\path{doi:10.2307/3621547}}.

\bibitem{teller89}
Paul Teller.
\newblock {\em A Modern Formal Logic Primer: Predicate Logic and Metatheory,
  Volume II}.
\newblock Prentice Hall, 1989.

\bibitem{trnkova71}
V\v{e}ra Trnkov\'a.
\newblock On descriptive classification of set-functors. {I}.
\newblock {\em Comment. Math. Univ. Carol.}, 12(1):143--174, 1971.

\bibitem{urbat17}
Henning Urbat.
\newblock {Finite Behaviours and Finitary Corecursion}.
\newblock In {\em 7th Conference on Algebra and Coalgebra in Computer Science
  (CALCO 2017)}, volume~72 of {\em LIPIcs}, pages 24:1--24:16. Schloss Dagstuhl
  -- Leibniz-Zentrum f{\"u}r Informatik, 2017.
\newblock \href {https://doi.org/10.4230/LIPIcs.CALCO.2017.24}
  {\path{doi:10.4230/LIPIcs.CALCO.2017.24}}.

\bibitem{wm24}
Thorsten Wi\ss{}mann and Stefan Milius.
\newblock Initial algebras unchained - {A} novel initial algebra construction
  formalized in {A}gda.
\newblock In {\em 9th Annual ACM/IEEE Symposium on Logic in Computer Science
  (LICS 2024)}. ACM, 2024.
\newblock \href {https://doi.org/10.1145/3661814.3662105}
  {\path{doi:10.1145/3661814.3662105}}.

\end{thebibliography}
\end{document}
